\documentclass[aps,preprint]{revtex4}%
\usepackage{amssymb}
\usepackage{amsfonts}
\usepackage{amsmath}
\usepackage{graphicx}%
\setcounter{MaxMatrixCols}{30}
%TCIDATA{OutputFilter=latex2.dll}
%TCIDATA{Version=5.50.0.2960}
%TCIDATA{CSTFile=revtex4.cst}
%TCIDATA{Created=Sunday, August 25, 2019 18:02:44}
%TCIDATA{LastRevised=Monday, May 22, 2023 07:01:03}
%TCIDATA{<META NAME="GraphicsSave" CONTENT="32">}
%TCIDATA{<META NAME="SaveForMode" CONTENT="1">}
%TCIDATA{BibliographyScheme=Manual}
%TCIDATA{<META NAME="DocumentShell" CONTENT="Articles\SW\REVTeX 4">}
%BeginMSIPreambleData
\providecommand{\U}[1]{\protect\rule{.1in}{.1in}}
%EndMSIPreambleData
\newtheorem{theorem}{Theorem}
\newtheorem{acknowledgement}[theorem]{Acknowledgement}
\newtheorem{algorithm}[theorem]{Algorithm}
\newtheorem{axiom}[theorem]{Axiom}
\newtheorem{claim}[theorem]{Claim}
\newtheorem{conclusion}[theorem]{Conclusion}
\newtheorem{condition}[theorem]{Condition}
\newtheorem{conjecture}[theorem]{Conjecture}
\newtheorem{corollary}[theorem]{Corollary}
\newtheorem{criterion}[theorem]{Criterion}
\newtheorem{definition}[theorem]{Definition}
\newtheorem{example}[theorem]{Example}
\newtheorem{exercise}[theorem]{Exercise}
\newtheorem{lemma}[theorem]{Lemma}
\newtheorem{notation}[theorem]{Notation}
\newtheorem{problem}[theorem]{Problem}
\newtheorem{proposition}[theorem]{Proposition}
\newtheorem{remark}[theorem]{Remark}
\newtheorem{solution}[theorem]{Solution}
\newtheorem{summary}[theorem]{Summary}
\newenvironment{proof}[1][Proof]{\noindent\textbf{#1.} }{\ \rule{0.5em}{0.5em}}
\begin{document}
\preprint{HEP/123-qed}
\title[Short title for running header]{REV\TeX4}
\author{First Author}
\affiliation{My Institution}
\author{Second Author}
\affiliation{My Institution}
\author{Third Author}
\affiliation{Other Institution}
\keywords{one two three}
\pacs{PACS number}

\begin{abstract}
Shell document for REV\TeX{} 4.

\end{abstract}
\volumeyear{year}
\volumenumber{number}
\issuenumber{number}
\eid{identifier}
\date[Date text]{date}
\received[Received text]{date}

\revised[Revised text]{date}

\accepted[Accepted text]{date}

\published[Published text]{date}

\startpage{101}
\endpage{102}
\maketitle
\tableofcontents


\section{Exterior Algebra and Mixed Exterior Algebra}

One defines a two sided ideal, $N\left(  E\right)  $ in $%
%TCIMACRO{\tbigotimes }%
%BeginExpansion
{\textstyle\bigotimes}
%EndExpansion
\left(  E\right)  $ generated by $\left\{  \left.  x\otimes x\right\vert x\in
E\right\}  ,$ where
\begin{equation}%
%TCIMACRO{\tbigotimes }%
%BeginExpansion
{\textstyle\bigotimes}
%EndExpansion
\left(  E\right)  =%
%TCIMACRO{\tbigoplus \limits_{0\leq p\leq n}}%
%BeginExpansion
{\textstyle\bigoplus\limits_{0\leq p\leq n}}
%EndExpansion%
%TCIMACRO{\tbigotimes \nolimits^{p}}%
%BeginExpansion
{\textstyle\bigotimes\nolimits^{p}}
%EndExpansion
E
\end{equation}
and elements of the ideal are given by
\begin{equation}
w_{1}\otimes\cdots\otimes w_{m}\otimes x_{1}\otimes x_{1}\otimes\cdots\otimes
x_{q}\otimes x_{q}\otimes y_{1}\otimes\cdots\otimes y_{s};\quad w_{j}%
,x_{k},y_{l}\in E;0\leq j\leq m\leq n,1\leq k\leq q,0\leq l\leq s.
\end{equation}
Greub shows that the ideal $N\left(  E\right)  $ is also generated by
\begin{gather}
w_{1}\otimes\cdots\otimes w_{m}\otimes\left(  x_{1}\otimes x_{2}+x_{2}\otimes
x_{1}\right)  \otimes\cdots\otimes\left(  x_{2q}\otimes x_{2q-1}%
+x_{2q-1}\otimes x_{2q}\right)  \otimes y_{1}\otimes\cdots\otimes y_{s}\\
w_{j},x_{k},y_{l}\in E;0\leq j\leq m\leq n,1\leq k\leq q,0\leq l\leq s.
\end{gather}
The ideal $N\left(  E\right)  $ is a subspace of $%
%TCIMACRO{\tbigotimes }%
%BeginExpansion
{\textstyle\bigotimes}
%EndExpansion
\left(  E\right)  $ and can also be expressed as
\begin{equation}
N^{p}\left(  E\right)  =\left\{  \left.  x_{j_{1}}\otimes\cdots\otimes
x_{j_{p}}\right\vert j_{k}=j_{l},k\neq l,\text{ for some }k.l\in\left\{
1,\cdots,p\right\}  \right\}  ;2\leq p\leq n,
\end{equation}
where $\left\{  \left.  x_{j_{1}}\otimes\cdots\otimes x_{j_{p}}\right\vert
0\leq p\leq n;1\leq j_{1}\,<\cdots<j_{p}\leq n\right\}  $ is a basis for $%
%TCIMACRO{\tbigotimes }%
%BeginExpansion
{\textstyle\bigotimes}
%EndExpansion
\left(  E\right)  .$

The ideal $N\left(  E\right)  $ defines an equivalence relationship between
elements of $%
%TCIMACRO{\tbigotimes }%
%BeginExpansion
{\textstyle\bigotimes}
%EndExpansion
\left(  E\right)  $ by
\begin{equation}
v\sim_{\wedge}w\Leftrightarrow v-w\in N\left(  E\right)
\end{equation}
that produces the eqivalence classes, $\left[  w\right]  _{\wedge}$, of
elements of $%
%TCIMACRO{\tbigotimes }%
%BeginExpansion
{\textstyle\bigotimes}
%EndExpansion
\left(  E\right)  $
\begin{equation}
\left[  w\right]  _{\wedge}=\left\{  \left.  w+N\left(  E\right)  \right\vert
\right\}  .
\end{equation}
The equivalence classes form a vector space i.e.
\begin{equation}
\alpha\left[  w_{1}\right]  _{\wedge}+\beta\left[  w_{2}\right]  _{\wedge
}=\left[  \alpha w_{1}+\beta w_{2}\right]  _{\wedge}%
\end{equation}
and thus is a subalgebra. The tensor product of equivalence classes gives
\begin{equation}
\left[  w_{1}\right]  _{\wedge}\otimes\left[  w_{2}\right]  _{\wedge}=\left[
w_{1}\otimes w_{2}\right]  _{\wedge}.
\end{equation}


The quotient space
\begin{equation}%
%TCIMACRO{\tbigwedge }%
%BeginExpansion
{\textstyle\bigwedge}
%EndExpansion
\left(  E\right)  =\frac{%
%TCIMACRO{\tbigotimes }%
%BeginExpansion
{\textstyle\bigotimes}
%EndExpansion
\left(  E\right)  }{\sim_{\wedge}}=\frac{%
%TCIMACRO{\tbigotimes }%
%BeginExpansion
{\textstyle\bigotimes}
%EndExpansion
\left(  E\right)  }{N\left(  E\right)  }%
\end{equation}
The antisymmetrizer $\mathcal{A}_{p}$ is a projection
\begin{equation}
\Pi_{A_{p}}\equiv\mathcal{A}_{p}:%
%TCIMACRO{\tbigotimes \nolimits^{p}}%
%BeginExpansion
{\textstyle\bigotimes\nolimits^{p}}
%EndExpansion
E\rightarrow%
%TCIMACRO{\tbigotimes \nolimits^{p}}%
%BeginExpansion
{\textstyle\bigotimes\nolimits^{p}}
%EndExpansion
E
\end{equation}
and can be related to $N\left(  E\right)  $ by observing
\begin{align}
N^{p}\left(  E\right)   &  =\ker\left\{  \Pi_{A_{p}}\right\}  \\
X^{p}\left(  E\right)   &  =\operatorname{Im}\left\{  \Pi_{A_{p}}\right\}
\end{align}%
\begin{equation}%
%TCIMACRO{\tbigotimes \nolimits^{p}}%
%BeginExpansion
{\textstyle\bigotimes\nolimits^{p}}
%EndExpansion
E=X^{p}\left(  E\right)
%TCIMACRO{\tbigoplus }%
%BeginExpansion
{\textstyle\bigoplus}
%EndExpansion
N^{p}\left(  E\right)  .
\end{equation}
The antisymmetrizer can be extended to all of $%
%TCIMACRO{\tbigotimes }%
%BeginExpansion
{\textstyle\bigotimes}
%EndExpansion
\left(  E\right)  $ by
\begin{gather}
\Pi_{A}=%
%TCIMACRO{\tsum \limits_{0\leq p\leq n}}%
%BeginExpansion
{\textstyle\sum\limits_{0\leq p\leq n}}
%EndExpansion
\mathcal{A}_{p}\equiv%
%TCIMACRO{\tsum \limits_{0\leq p\leq n}}%
%BeginExpansion
{\textstyle\sum\limits_{0\leq p\leq n}}
%EndExpansion
\Pi_{A_{p}}\\
\Pi_{A_{p}}\circ\Pi_{A_{q}}=\Pi_{A_{p}}\delta_{pq}\\
\Pi_{A}^{2}=%
%TCIMACRO{\tsum \limits_{0\leq p\leq n}}%
%BeginExpansion
{\textstyle\sum\limits_{0\leq p\leq n}}
%EndExpansion
\Pi_{A_{p}}\circ%
%TCIMACRO{\tsum \limits_{0\leq q\leq n}}%
%BeginExpansion
{\textstyle\sum\limits_{0\leq q\leq n}}
%EndExpansion
\Pi_{A_{q}}=%
%TCIMACRO{\tsum \limits_{0\leq p,q\leq n}}%
%BeginExpansion
{\textstyle\sum\limits_{0\leq p,q\leq n}}
%EndExpansion
\Pi_{A_{p}}\circ\Pi_{A_{q}}=%
%TCIMACRO{\tsum \limits_{0\leq p\leq n}}%
%BeginExpansion
{\textstyle\sum\limits_{0\leq p\leq n}}
%EndExpansion
\Pi_{A_{p}}=\Pi_{A}%
\end{gather}
and%
\begin{align}
N\left(  E\right)   &  =%
%TCIMACRO{\tbigoplus \limits_{0\leq p\leq n}}%
%BeginExpansion
{\textstyle\bigoplus\limits_{0\leq p\leq n}}
%EndExpansion
N^{p}\left(  E\right)  \\
X\left(  E\right)   &  =%
%TCIMACRO{\tbigoplus \limits_{0\leq p\leq n}}%
%BeginExpansion
{\textstyle\bigoplus\limits_{0\leq p\leq n}}
%EndExpansion
X^{p}\left(  E\right)
\end{align}
showing that the tensor algebra can be decomposed as
\begin{equation}%
%TCIMACRO{\tbigotimes }%
%BeginExpansion
{\textstyle\bigotimes}
%EndExpansion
\left(  E\right)  =X\left(  E\right)
%TCIMACRO{\tbigoplus }%
%BeginExpansion
{\textstyle\bigoplus}
%EndExpansion
N\left(  E\right)  .
\end{equation}
The inner product of antisymmetrized tensors is
\begin{equation}
\left\langle \Pi_{A_{p}}^{d}\left(  y_{1}^{d}\otimes\cdots\otimes y_{p}%
^{d}\right)  ,\Pi_{A_{p}}\left(  x_{1}\otimes\cdots\otimes x_{p}\right)
\right\rangle =\frac{1}{p!}\det\left\{  \left\langle y_{j}^{d},x_{k}%
\right\rangle \right\}  ,
\end{equation}
The antisymmetrizer and its orthogonal compliment have the following
properties
\begin{align}
x &  =\Pi_{A}\left(  x\right)  +\Pi_{A}^{\bot}\left(  x\right)  \\
y &  =\Pi_{A}\left(  y\right)  +\Pi_{A}^{\bot}\left(  y\right)  \\
x\otimes y &  =\Pi_{A}\left(  x\otimes y\right)  +\Pi_{A}^{\bot}\left(
x\otimes y\right)  \\
x\otimes y &  =\left(  \Pi_{A}\left(  x\right)  +\Pi_{A}^{\bot}\left(
x\right)  \right)  \otimes\left(  \Pi_{A}\left(  y\right)  +\Pi_{A}^{\bot
}\left(  y\right)  \right)  \\
&  =\Pi_{A}\left(  x\right)  \otimes\Pi_{A}\left(  y\right)  +\Pi_{A}\left(
x\right)  \otimes\Pi_{A}^{\bot}\left(  y\right)  +\Pi_{A}^{\bot}\left(
x\right)  \otimes\Pi_{A}\left(  y\right)  +\Pi_{A}^{\bot}\left(  x\right)
\otimes\Pi_{A}^{\bot}\left(  y\right)  \\
&  =\left(  \Pi_{A}\otimes\Pi_{A}\right)  \left(  x\otimes y\right)  +\left(
\Pi_{A}\otimes\Pi_{A}^{\bot}\right)  \left(  x\otimes y\right)  +\left(
\Pi_{A}^{\bot}\otimes\Pi_{A}\right)  \left(  x\otimes y\right)  +\left(
\Pi_{A}^{\bot}\otimes\Pi_{A}^{\bot}\right)  \left(  x\otimes y\right)  \\
.\Pi_{A}\left(  x\otimes y\right)   &  =\Pi_{A}\left(  \Pi_{A}\left(
x\right)  \otimes\Pi_{A}\left(  y\right)  +\Pi_{A}\left(  x\right)  \otimes
\Pi_{A}^{\bot}\left(  y\right)  +\Pi_{A}^{\bot}\left(  x\right)  \otimes
\Pi_{A}\left(  y\right)  +\Pi_{A}^{\bot}\left(  x\right)  \otimes\Pi_{A}%
^{\bot}\left(  y\right)  \right)  .
\end{align}
These relationships produced the very impotant identity%
\begin{equation}
\Pi_{A}\left(  x\otimes y\right)  =\Pi_{A}\left(  \Pi_{A}x\otimes\Pi
_{A}y\right)  ,x,y\in%
%TCIMACRO{\tbigotimes }%
%BeginExpansion
{\textstyle\bigotimes}
%EndExpansion
\left(  E\right)  .
\end{equation}
The mapping onto the quotient space is
\begin{equation}
\Pi:%
%TCIMACRO{\tbigotimes }%
%BeginExpansion
{\textstyle\bigotimes}
%EndExpansion
\left(  E\right)  \rightarrow\left.
%TCIMACRO{\tbigotimes }%
%BeginExpansion
{\textstyle\bigotimes}
%EndExpansion
\left(  E\right)  \right/  N\left(  E\right)
\end{equation}


One has the commutative diagram%
\begin{align}
&
%TCIMACRO{\tbigotimes }%
%BeginExpansion
{\textstyle\bigotimes}
%EndExpansion
\left(  E\right)  \longrightarrow\longrightarrow\longrightarrow\overset{\pi
_{X}}{\longrightarrow}\longrightarrow\longrightarrow X\left(  E\right)
\subset%
%TCIMACRO{\tbigotimes }%
%BeginExpansion
{\textstyle\bigotimes}
%EndExpansion
\left(  E\right) \\
&  \qquad\downarrow\qquad\qquad\qquad\swarrow\\
&  \,\left.  \pi\right.  \,\downarrow\qquad\qquad\swarrow\left.  \rho\right.
\\
&  \qquad\downarrow\qquad\swarrow\\
&  \qquad\left.
%TCIMACRO{\tbigotimes }%
%BeginExpansion
{\textstyle\bigotimes}
%EndExpansion
\left(  E\right)  \right/  N\left(  E\right)  ,
\end{align}
where%
\begin{align}
&  \Pi:%
%TCIMACRO{\tbigotimes }%
%BeginExpansion
{\textstyle\bigotimes}
%EndExpansion
\left(  E\right)  =\sum_{0\leq p\leq n}%
%TCIMACRO{\tbigotimes \nolimits^{p}}%
%BeginExpansion
{\textstyle\bigotimes\nolimits^{p}}
%EndExpansion
E\rightarrow\left.
%TCIMACRO{\tbigotimes }%
%BeginExpansion
{\textstyle\bigotimes}
%EndExpansion
\left(  E\right)  \right/  N\left(  E\right)  =\sum_{0\leq p\leq n}\Pi\left(
%TCIMACRO{\tbigotimes \nolimits^{p}}%
%BeginExpansion
{\textstyle\bigotimes\nolimits^{p}}
%EndExpansion
E\right) \\
&  \Pi_{X}:%
%TCIMACRO{\tbigotimes }%
%BeginExpansion
{\textstyle\bigotimes}
%EndExpansion
\left(  E\right)  \rightarrow%
%TCIMACRO{\tbigotimes }%
%BeginExpansion
{\textstyle\bigotimes}
%EndExpansion
\left(  E\right) \\
&  \left.  \Pi\right\vert _{X\left(  E\right)  }=\rho:X\left(  E\right)
\rightarrow\left.
%TCIMACRO{\tbigotimes }%
%BeginExpansion
{\textstyle\bigotimes}
%EndExpansion
\left(  E\right)  \right/  N\left(  E\right)  .
\end{align}
The sum is the same as the direct sum as the subspaces $%
%TCIMACRO{\tbigotimes \nolimits^{p}}%
%BeginExpansion
{\textstyle\bigotimes\nolimits^{p}}
%EndExpansion
E$ are disjoint subspaces of $%
%TCIMACRO{\tbigotimes }%
%BeginExpansion
{\textstyle\bigotimes}
%EndExpansion
\left(  E\right)  .$

Since $N\left(  E\right)  $ is an ideal in $%
%TCIMACRO{\tbigotimes }%
%BeginExpansion
{\textstyle\bigotimes}
%EndExpansion
\left(  E\right)  $ then
\begin{equation}
\Pi\left(  a\right)  \wedge_{G}\Pi\left(  b\right)  =\Pi\left(  a\otimes
b\right)  ;a,b\in%
%TCIMACRO{\tbigotimes }%
%BeginExpansion
{\textstyle\bigotimes}
%EndExpansion
\left(  E\right)
\end{equation}
and
\begin{equation}
\left\langle \Pi^{d}\left(  y_{1}^{d}\otimes\cdots\otimes y_{p}^{d}\right)
,\Pi\left(  x_{1}\otimes\cdots\otimes x_{p}\right)  \right\rangle
=p!\left\langle \Pi_{A_{p}}^{d}\left(  y_{1}^{d}\otimes\cdots\otimes y_{p}%
^{d}\right)  ,\Pi_{A_{p}}\left(  x_{1}\otimes\cdots\otimes x_{p}\right)
\right\rangle =\det\left\{  \left\langle y_{j}^{d},x_{k}\right\rangle
\right\}  .
\end{equation}
If $a,b\in E$ then
\begin{equation}
\Pi_{X}\left(  a\otimes b\right)  =a\otimes b-b\otimes a.
\end{equation}


\begin{notation}
Greub defines exterior product "$\wedge_{G}$" as
\begin{equation}
\wedge_{G}\left(  x_{j_{1}}\otimes\cdots\otimes x_{j_{p}}\right)  =x_{j_{1}%
}\wedge_{G}\cdots\wedge_{G}x_{j_{p}}=\sum_{\sigma\in S_{p}}\left(  -1\right)
^{\pi\left(  \sigma\right)  }x_{j_{\sigma\left(  1\right)  }}\otimes
\cdots\otimes x_{j_{\sigma\left(  p\right)  }},
\end{equation}
while he defines another product "$\cap_{G}$" as
\begin{equation}
\cap_{G}\left(  x_{j_{1}}\otimes\cdots\otimes x_{j_{p}}\right)  =x_{j_{1}}%
\cap\cdots\cap x_{j_{p}}=\frac{1}{p!}\sum_{\sigma\in S_{p}}\left(  -1\right)
^{\pi\left(  \sigma\right)  }x_{j_{\sigma\left(  1\right)  }}\otimes
\cdots\otimes x_{j_{\sigma\left(  p\right)  }}.
\end{equation}


\begin{example}%
\begin{equation}
\wedge_{G}\left(  x_{1}\otimes x_{2}\right)  =x_{1}\wedge_{G}x_{2}%
=x_{1}\otimes x_{2}-x_{2}\otimes x_{1},
\end{equation}
$\lceil$%
\begin{align}
\wedge_{G}^{2}\left(  x_{1}\otimes x_{2}\right)   &  =\wedge_{G}\left(
x_{1}\otimes x_{2}-x_{2}\otimes x_{1}\right)  =\wedge_{G}\left(  x_{1}\otimes
x_{2}\right)  -\wedge_{G}\left(  x_{2}\otimes x_{1}\right) \\
&  =\left(  x_{1}\otimes x_{2}-x_{2}\otimes x_{1}\right)  -\left(
x_{2}\otimes x_{1}-x_{1}\otimes x_{2}\right) \\
&  =2\left(  x_{1}\otimes x_{2}-x_{2}\otimes x_{1}\right)  ,
\end{align}
$\rfloor$

and
\begin{equation}
\cap_{G}\left(  x_{1}\otimes x_{2}\right)  =x_{1}\cap_{G}x_{2}=\frac{1}%
{2}\left(  x_{1}\otimes x_{2}-x_{2}\otimes x_{1}\right)  .
\end{equation}
Multiple products are defined by
\begin{align}
&  x_{j_{1}}\wedge_{G}\cdots\wedge_{G}x_{j_{p}}\wedge_{G}x_{j_{p+1}}\wedge
_{G}\cdots\wedge_{G}x_{j_{p+q}}=\left(  x_{j_{1}}\wedge_{G}\cdots\wedge
_{G}x_{j_{p}}\right)  \wedge_{G}\left(  x_{j_{p+1}}\wedge_{G}\cdots\wedge
_{G}x_{j_{p+q}}\right) \\
&  =\wedge_{G}\left(  \left(  \wedge_{G}\left(  x_{j_{1}}\otimes\cdots\otimes
x_{j_{p}}\right)  \right)  \otimes\left(  \wedge_{G}\left(  x_{j_{p+1}}%
\otimes\cdots\otimes x_{j_{p+q}}\right)  \right)  \right)  =\left(  \wedge
_{G}\left(  x_{j_{1}}\otimes\cdots\otimes x_{j_{p}}\right)  \right)
\wedge_{G}\left(  \wedge_{G}\left(  x_{j_{p+1}}\otimes\cdots\otimes
x_{j_{p+q}}\right)  \right)
\end{align}
and
\begin{align}
&  x_{j_{1}}\cap_{G}\cdots\cap_{G}x_{j_{p}}\cap_{G}x_{j_{p+1}}\cap_{G}%
\cdots\cap_{G}x_{j_{p+q}}=\left(  x_{j_{1}}\cap_{G}\cdots\cap_{G}x_{j_{p}%
}\right)  \cap_{G}\left(  x_{j_{p+1}}\cap_{G}\cdots\cap_{G}x_{j_{p+q}}\right)
\\
&  =\cap_{G}\left(  \left(  \cap_{G}\left(  x_{j_{1}}\otimes\cdots\otimes
x_{j_{p}}\right)  \right)  \otimes\left(  \cap_{G}\left(  x_{j_{p+1}}%
\otimes\cdots\otimes x_{j_{p+q}}\right)  \right)  \right)  =\left(  \cap
_{G}\left(  x_{j_{1}}\otimes\cdots\otimes x_{j_{p}}\right)  \right)  \cap
_{G}\left(  \cap_{G}\left(  x_{j_{p+1}}\otimes\cdots\otimes x_{j_{p+q}%
}\right)  \right)  .
\end{align}
In particular
\begin{align}
&  x_{j_{1}}\wedge_{G}\cdots\wedge_{G}x_{j_{p}}\wedge_{G}x_{j_{p+1}}=\left(
x_{j_{1}}\wedge_{G}\cdots\wedge_{G}x_{j_{p}}\right)  \wedge_{G}\left(
x_{j_{p+1}}\right) \\
&  =\wedge_{G}\left(  \left(  \wedge_{G}\left(  x_{j_{1}}\otimes\cdots\otimes
x_{j_{p}}\right)  \right)  \otimes\left(  \wedge_{G}\left(  x_{j_{p+1}%
}\right)  \right)  \right)  =\left(  \wedge_{G}\left(  x_{j_{1}}\otimes
\cdots\otimes x_{j_{p}}\right)  \right)  \wedge_{G}\left(  \wedge_{G}\left(
x_{j_{p+1}}\right)  \right)
\end{align}
and
\begin{align}
&  x_{j_{1}}\wedge_{G}\cdots\wedge_{G}x_{j_{p}}\wedge_{G}x_{j_{p+1}}=\left(
x_{j_{1}}\wedge_{G}\cdots\wedge_{G}x_{j_{p}}\right)  \wedge_{G}\left(
x_{j_{p+1}}\right) \\
&  =\wedge_{G}\left(  \left(  \wedge_{G}\left(  x_{j_{1}}\otimes\cdots\otimes
x_{j_{p}}\right)  \right)  \otimes\left(  \wedge_{G}\left(  x_{j_{p+1}%
}\right)  \right)  \right)  =\left(  \wedge_{G}\left(  x_{j_{1}}\otimes
\cdots\otimes x_{j_{p}}\right)  \right)  \wedge_{G}\left(  \wedge_{G}\left(
x_{j_{p+1}}\right)  \right)
\end{align}

\end{example}

\begin{example}
Consider
\begin{equation}
\cap_{G}\left(  \left(  x_{1}\cap_{G}x_{2}\right)  \otimes\left(  x_{3}%
\cap_{G}x_{4}\right)  \right)  =\left(  x_{1}\cap_{G}x_{2}\right)  \cap
_{G}\left(  x_{3}\cap_{G}x_{4}\right)
\end{equation}
the product of the two individual factors gives%
\begin{align}
&  \left(  \frac{1}{2!}\sum_{\sigma\in S_{2}}\left(  -1\right)  ^{\pi\left(
\sigma\right)  }x_{\sigma\left(  1\right)  }\otimes x_{\sigma\left(  2\right)
}\right)  \cap_{G}\left(  \frac{1}{2!}\sum_{\mu\in S_{2}}\left(  -1\right)
^{\pi\left(  \mu\right)  }x_{\mu\left(  3\right)  }\otimes x_{\mu\left(
4\right)  }\right) \\
&  =\frac{1}{4}\left(  x_{1}\otimes x_{2}-x_{2}\otimes x_{1}\right)  \cap
_{G}\left(  x_{3}\otimes x_{4}-x_{4}\otimes x_{3}\right) \\
&  =\frac{1}{4}\left(  \left(  x_{1}\otimes x_{2}\right)  \cap_{G}\left(
x_{3}\otimes x_{4}\right)  -\left(  x_{1}\otimes x_{2}\right)  \cap_{G}\left(
x_{4}\otimes x_{3}\right)  -\left(  x_{2}\otimes x_{1}\right)  \cap_{G}\left(
x_{3}\otimes x_{4}\right)  +\left(  x_{2}\otimes x_{1}\right)  \cap_{G}\left(
x_{4}\otimes x_{3}\right)  \right)
\end{align}
and the final product gives
\begin{gather}
\frac{1}{4}\left(  \mathcal{A}_{4}\left(  x_{1}\otimes x_{2}\otimes
x_{3}\otimes x_{4}\right)  -\mathcal{A}_{4}\left(  x_{1}\otimes x_{2}\otimes
x_{4}\otimes x_{3}\right)  -\mathcal{A}_{4}\left(  x_{2}\otimes x_{1}\otimes
x_{3}\otimes x_{4}\right)  +\mathcal{A}_{4}\left(  x_{2}\otimes x_{1}\otimes
x_{4}\otimes x_{3}\right)  \right) \\
=\mathcal{A}_{4}\left(  x_{1}\otimes x_{2}\otimes x_{3}\otimes x_{4}\right)
\end{gather}
i.e.
\begin{align}
&  \frac{1}{4!}\sum_{\sigma\in S_{4}}\left(  -1\right)  ^{\pi\left(
\sigma\right)  }\frac{1}{4}\left(  x_{\sigma\left(  1\right)  }\otimes
x_{\sigma\left(  2\right)  }\otimes x_{\sigma\left(  3\right)  }\otimes
x_{\sigma\left(  4\right)  }-x_{\sigma\left(  1\right)  }\otimes
x_{\sigma\left(  2\right)  }\otimes x_{\sigma\left(  4\right)  }\otimes
x_{\sigma\left(  3\right)  }\right. \\
&  =-\left.  x_{\sigma\left(  2\right)  }\otimes x_{\sigma\left(  1\right)
}\otimes x_{\sigma\left(  3\right)  }\otimes x_{\sigma\left(  4\right)
}+x_{\sigma\left(  2\right)  }\otimes x_{\sigma\left(  1\right)  }\otimes
x_{\sigma\left(  4\right)  }\otimes x_{\sigma\left(  3\right)  }\right) \\
&  =\frac{1}{4!}\sum_{\sigma\in S_{4}}\left(  -1\right)  ^{\pi\left(
\sigma\right)  }x_{\sigma\left(  1\right)  }\otimes x_{\sigma\left(  2\right)
}\otimes x_{\sigma\left(  3\right)  }\otimes x_{\sigma\left(  4\right)  }.
\end{align}
The Greub product of individual factors
\begin{equation}
\wedge_{G}\left(  \left(  x_{1}\wedge_{G}x_{2}\right)  \otimes\left(
x_{3}\wedge_{G}x_{4}\right)  \right)  =\left(  x_{1}\wedge_{G}x_{2}\right)
\wedge_{G}\left(  x_{3}\wedge_{G}x_{4}\right)
\end{equation}
gives
\begin{align}
&  \left(  \sum_{\sigma\in S_{2}}\left(  -1\right)  ^{\pi\left(
\sigma\right)  }x_{\sigma\left(  1\right)  }\otimes x_{\sigma\left(  2\right)
}\right)  \wedge_{G}\left(  \sum_{\mu\in S_{2}}\left(  -1\right)  ^{\pi\left(
\mu\right)  }x_{\mu\left(  3\right)  }\otimes x_{\mu\left(  4\right)  }\right)
\\
&  =\left(  x_{1}\otimes x_{2}-x_{2}\otimes x_{1}\right)  \wedge_{G}\left(
x_{3}\otimes x_{4}-x_{4}\otimes x_{3}\right) \\
&  =\left(  x_{1}\otimes x_{2}\right)  \wedge_{G}\left(  x_{3}\otimes
x_{4}\right)  -\left(  x_{1}\otimes x_{2}\right)  \wedge_{G}\left(
x_{4}\otimes x_{3}\right)  -\left(  x_{2}\otimes x_{1}\right)  \wedge
_{G}\left(  x_{3}\otimes x_{4}\right)  +\left(  x_{2}\otimes x_{1}\right)
\wedge_{G}\left(  x_{4}\otimes x_{3}\right)
\end{align}
and the final product gives%
\begin{align}
&  \sum_{\sigma\in S_{4}}\left(  -1\right)  ^{\pi\left(  \sigma\right)
}\left(  x_{\sigma\left(  1\right)  }\otimes x_{\sigma\left(  2\right)
}\otimes x_{\sigma\left(  3\right)  }\otimes x_{\sigma\left(  4\right)
}-x_{\sigma\left(  1\right)  }\otimes x_{\sigma\left(  2\right)  }\otimes
x_{\sigma\left(  4\right)  }\otimes x_{\sigma\left(  3\right)  }\right. \\
&  =-\left.  x_{\sigma\left(  2\right)  }\otimes x_{\sigma\left(  1\right)
}\otimes x_{\sigma\left(  3\right)  }\otimes x_{\sigma\left(  4\right)
}+x_{\sigma\left(  2\right)  }\otimes x_{\sigma\left(  1\right)  }\otimes
x_{\sigma\left(  4\right)  }\otimes x_{\sigma\left(  3\right)  }\right) \\
&  =4\sum_{\sigma\in S_{4}}\left(  -1\right)  ^{\pi\left(  \sigma\right)
}x_{\sigma\left(  1\right)  }\otimes x_{\sigma\left(  2\right)  }\otimes
x_{\sigma\left(  3\right)  }\otimes x_{\sigma\left(  4\right)  }.
\end{align}
However
\begin{equation}
4\sum_{\sigma\in S_{4}}\left(  -1\right)  ^{\pi\left(  \sigma\right)
}x_{\sigma\left(  1\right)  }\otimes x_{\sigma\left(  2\right)  }\otimes
x_{\sigma\left(  3\right)  }\otimes x_{\sigma\left(  4\right)  }\neq
x_{1}\wedge_{G}x_{2}\wedge_{G}x_{3}\wedge_{G}x_{4}.
\end{equation}

\end{example}
\end{notation}

\begin{remark}
The products are associative i.e.
\begin{equation}
x_{1}\cap_{G}\left(  x_{2}\cap_{G}x_{3}\right)  =\left(  x_{1}\cap_{G}%
x_{2}\right)  \cap_{G}x_{3}%
\end{equation}
which is defined by
\begin{equation}
\mathcal{A}_{3}\left(  x_{1}\otimes\mathcal{A}_{2}\left(  x_{2}\otimes
x_{3}\right)  \right)  =\mathcal{A}_{3}\left(  \mathcal{A}_{2}\left(
x_{1}\otimes x_{2}\right)  \otimes x_{3}\right)  .
\end{equation}
one can show that
\begin{align}
&  \mathcal{A}_{p+q}\left(  \mathcal{A}_{p}\left(  x_{j_{1}}\otimes
\cdots\otimes x_{j_{p}}\right)  \otimes\mathcal{A}_{q}\left(  x_{j_{p+1}%
}\otimes\cdots\otimes x_{j_{p+q}}\right)  \right) \\
&  =\mathcal{A}_{p+q}\left(  \frac{1}{p!}\sum_{\sigma\in S_{p}}\left(
-1\right)  ^{\pi\left(  \sigma\right)  }x_{j_{\sigma\left(  1\right)  }%
}\otimes\cdots\otimes x_{j_{\sigma\left(  p\right)  }}\otimes\frac{1}{q!}%
\sum_{\mu\in S_{q}}\left(  -1\right)  ^{\pi\left(  \mu\right)  }%
x_{j_{\mu\left(  p+1\right)  }}\otimes\cdots\otimes x_{j_{\mu\left(
P+q\right)  }}\right) \\
&  =\frac{1}{\left(  p+q\right)  !}\sum_{\lambda\in S_{p+q}}\left(  -1\right)
^{\pi\left(  \lambda\right)  }x_{j_{\lambda\left(  1\right)  }}\otimes
\cdots\otimes x_{j_{\lambda\left(  p\right)  }}\otimes x_{j_{\lambda\left(
p+1\right)  }}\otimes\cdots\otimes x_{j_{\lambda\left(  P+q\right)  }}.
\end{align}
Thus the equality
\begin{equation}
\mathcal{A}_{3}\left(  x_{1}\otimes\mathcal{A}_{2}\left(  x_{2}\otimes
x_{3}\right)  \right)  =\mathcal{A}_{3}\left(  \mathcal{A}_{2}\left(
x_{1}\otimes x_{2}\right)  \otimes x_{3}\right)
\end{equation}
is valid and the $\cap_{G}$ product is associative.
\end{remark}

\paragraph{Exterior Powers}

He defines the $k^{th}$ exterior power to be
\[
u^{k}=\frac{1}{k!}u\wedge_{G}\cdots\wedge_{G}u
\]
$\lceil$%
\begin{align}
u^{k}\wedge_{G}u^{l}  &  =\frac{1}{k!}\overleftarrow{u\wedge_{G}%
}\overset{k}{\cdots}\overrightarrow{\wedge_{G}u}\wedge_{G}\frac{1}%
{l!}\overleftarrow{u\wedge_{G}}\overset{l}{\cdots}\overrightarrow{\wedge_{G}%
u}=\frac{1}{k!}\frac{1}{l!}\overleftarrow{u\wedge_{G}}\overset{k+l}{\cdots
}\overrightarrow{\wedge_{G}u}\\
u^{k+l}  &  =\frac{1}{\left(  k+l\right)  !}\overleftarrow{u\wedge_{G}%
}\overset{k+l}{\cdots}\overrightarrow{\wedge_{G}u}\\
u^{k}\wedge_{G}u^{l}  &  =\alpha u^{k+l}\\
\frac{1}{k!}\frac{1}{l!}\overleftarrow{u\wedge_{G}}\overset{k+l}{\cdots
}\overrightarrow{\wedge_{G}u}  &  =\alpha\frac{1}{\left(  k+l\right)
!}\overleftarrow{u\wedge_{G}}\overset{k+l}{\cdots}\overrightarrow{\wedge_{G}%
u}\\
\alpha &  =\frac{\left(  k+l\right)  !}{k!l!}=\frac{\left(  k+l\right)
!}{\left(  k+l-l\right)  !l!}=\binom{k+l}{l}\\
u^{k}\wedge_{G}u^{l}  &  =\binom{k+l}{l}u^{k+l}%
\end{align}
$\rfloor$

\begin{example}%
\begin{align*}
&  \left(  x_{1}\wedge_{G}x_{2}+x_{3}\wedge_{G}x_{4}+x_{5}\wedge_{G}%
x_{6}+x_{7}\wedge_{G}x_{8}\right)  \wedge_{G}\left(  x_{1}\wedge_{G}%
x_{2}+x_{3}\wedge_{G}x_{4}+x_{5}\wedge_{G}x_{6}+x_{7}\wedge_{G}x_{8}\right) \\
&  =x_{1}\wedge_{G}x_{2}\wedge_{G}x_{3}\wedge_{G}x_{4}+x_{1}\wedge_{G}%
x_{2}\wedge_{G}x_{5}\wedge_{G}x_{6}+x_{1}\wedge_{G}x_{2}\wedge_{G}x_{7}%
\wedge_{G}x_{8}\\
&  +x_{3}\wedge_{G}x_{4}\wedge_{G}x_{1}\wedge_{G}x_{2}+x_{3}\wedge_{G}%
x_{4}\wedge_{G}x_{5}\wedge_{G}x_{6}+x_{3}\wedge_{G}x_{4}\wedge_{G}x_{7}%
\wedge_{G}x_{8}\\
&  +x_{5}\wedge_{G}x_{6}\wedge_{G}x_{1}\wedge_{G}x_{2}+x_{5}\wedge_{G}%
x_{6}\wedge_{G}x_{3}\wedge_{G}x_{4}+x_{5}\wedge_{G}x_{6}\wedge_{G}x_{7}%
\wedge_{G}x_{8}\\
&  +x_{7}\wedge_{G}x_{8}\wedge_{G}x_{1}\wedge_{G}x_{2}+x_{7}\wedge_{G}%
x_{8}\wedge_{G}x_{3}\wedge_{G}x_{4}+x_{7}\wedge_{G}x_{8}\wedge_{G}x_{5}%
\wedge_{G}x_{6}\\
&  =2\left(  x_{1}\wedge_{G}x_{2}\wedge_{G}x_{3}\wedge_{G}x_{4}+x_{1}%
\wedge_{G}x_{2}\wedge_{G}x_{5}\wedge_{G}x_{6}+x_{1}\wedge_{G}x_{2}\wedge
_{G}x_{7}\wedge_{G}x_{8}\right. \\
&  +\left.  x_{3}\wedge_{G}x_{4}\wedge_{G}x_{5}\wedge_{G}x_{6}+x_{3}\wedge
_{G}x_{4}\wedge_{G}x_{7}\wedge_{G}x_{8}+x_{5}\wedge_{G}x_{6}\wedge_{G}%
x_{7}\wedge_{G}x_{8}\right)  ,
\end{align*}
which gives%
\begin{align}
&  \left(  x_{1}\wedge_{G}x_{2}+x_{3}\wedge_{G}x_{4}+x_{5}\wedge_{G}%
x_{6}+x_{7}\wedge_{G}x_{8}\right)  ^{2}\\
&  \overset{\Delta}{=}\frac{1}{2}\left(  x_{1}\wedge_{G}x_{2}+x_{3}\wedge
_{G}x_{4}+x_{5}\wedge_{G}x_{6}+x_{7}\wedge_{G}x_{8}\right)  \wedge_{G}\left(
x_{1}\wedge_{G}x_{2}+x_{3}\wedge_{G}x_{4}+x_{5}\wedge_{G}x_{6}+x_{7}\wedge
_{G}x_{8}\right) \\
&  =\left(  x_{1}\wedge_{G}x_{2}\wedge_{G}x_{3}\wedge_{G}x_{4}+x_{1}\wedge
_{G}x_{2}\wedge_{G}x_{5}\wedge_{G}x_{6}+x_{1}\wedge_{G}x_{2}\wedge_{G}%
x_{7}\wedge_{G}x_{8}\right. \\
&  +\left.  x_{3}\wedge_{G}x_{4}\wedge_{G}x_{5}\wedge_{G}x_{6}+x_{3}\wedge
_{G}x_{4}\wedge_{G}x_{7}\wedge_{G}x_{8}+x_{5}\wedge_{G}x_{6}\wedge_{G}%
x_{7}\wedge_{G}x_{8}\right)  .
\end{align}

\end{example}

\begin{notation}
I define "$\wedge_{W}$" as his $\cap_{G}$
\begin{equation}
x\wedge_{W}y=\frac{1}{2}\left(  x_{1}\otimes x_{2}-x_{2}\otimes x_{1}\right)
=\mathcal{A}_{2}\left(  x_{1}\otimes x_{2}\right)  =\frac{1}{2}x_{1}\wedge
_{G}x_{2}.
\end{equation}
A ramification of this is that normailzed vectors $\left\vert x_{1}%
x_{2}\right\rangle $ are given in Greubs notation as
\begin{equation}
\left\vert x_{1}x_{2}\right\rangle =\frac{1}{\sqrt{2}}\left\vert x_{1}%
\wedge_{G}x_{2}\right\rangle =\frac{1}{\sqrt{2}}\left(  x_{1}\otimes
x_{2}-x_{2}\otimes x_{1}\right)
\end{equation}
in contrast to the definition in my notation
\begin{equation}
\left\vert x_{1}x_{2}\right\rangle =\sqrt{2}\left\vert x_{1}\wedge_{W}%
x_{2}\right\rangle =\frac{1}{\sqrt{2}}\left(  x_{1}\otimes x_{2}-x_{2}\otimes
x_{1}\right)  .
\end{equation}
The projectors are thus given by
\begin{equation}
\left\vert x_{1}x_{2}\right\rangle \left\langle x_{1}x_{2}\right\vert
=\frac{1}{2}\left\vert x_{1}\wedge_{G}x_{2}\right\rangle \left\langle
x_{1}\wedge_{G}x_{2}\right\vert =2\left\vert x_{1}\wedge_{W}x_{2}\right\rangle
\left\langle x_{1}\wedge_{W}x_{2}\right\vert
\end{equation}
and products of projectors as
\begin{equation}
\left\vert x_{1}\right\rangle \left\langle x_{1}\right\vert \boxtimes
_{G}\left\vert x_{2}\right\rangle \left\langle x_{2}\right\vert =\left\vert
x_{1}\wedge_{G}x_{2}\right\rangle \left\langle x_{1}\wedge_{G}x_{2}\right\vert
=2\left\vert x_{1}x_{2}\right\rangle \left\langle x_{1}x_{2}\right\vert
\end{equation}
and
\begin{equation}
\left\vert x_{1}\right\rangle \left\langle x_{1}\right\vert \boxtimes
_{W}\left\vert x_{2}\right\rangle \left\langle x_{2}\right\vert =\left\vert
x_{1}\wedge_{W}x_{2}\right\rangle \left\langle x_{1}\wedge_{W}x_{2}\right\vert
=\frac{1}{2}\left\vert x_{1}x_{2}\right\rangle \left\langle x_{1}%
x_{2}\right\vert .
\end{equation}

\end{notation}

\paragraph{Determinants}%

\begin{equation}
\left\langle \mathcal{A}_{p}^{d}y_{1}^{d}\otimes\cdots\otimes y_{p}%
^{d},\mathcal{A}_{p}x_{1}\otimes\cdots\otimes x_{p}\right\rangle =\left\langle
y_{1}^{d}\wedge_{W}\cdots\wedge_{W}y_{p}^{d},x_{1}\wedge_{W}\cdots\wedge
_{W}x_{p}\right\rangle =\frac{1}{p!}\det\left\{  \left\langle y_{j}^{d}%
,x_{k}\right\rangle \right\}  ,
\end{equation}
while
\begin{equation}
\left\langle y_{1}^{d}\wedge_{G}\cdots\wedge_{G}y_{p}^{d},x_{1}\wedge
_{G}\cdots\wedge_{G}x_{p}\right\rangle =\det\left\{  \left\langle y_{j}%
^{d},x_{k}\right\rangle \right\}  .
\end{equation}


\paragraph{Products of Projectors onto $\mathcal{H}^{1}$}

\begin{notation}
In the following we do not distinguish between the mixed exterior algebra $%
%TCIMACRO{\tbigwedge }%
%BeginExpansion
{\textstyle\bigwedge}
%EndExpansion
\left(  \mathcal{H}_{\mathcal{F}}\right)  \otimes%
%TCIMACRO{\tbigwedge }%
%BeginExpansion
{\textstyle\bigwedge}
%EndExpansion
\left(  \mathcal{H}_{\mathcal{F}}^{d}\right)  $ and its isomorphic image
\begin{equation}
\mathcal{B}_{00}\left(  \mathcal{H}_{\mathcal{F}}\right)  =\Xi_{\mathcal{H}%
_{\mathcal{F}}\mathcal{H}_{\mathcal{F}}^{d}}\left(  \overset{}{%
%TCIMACRO{\tbigwedge }%
%BeginExpansion
{\textstyle\bigwedge}
%EndExpansion
\left(  \mathcal{H}_{\mathcal{F}}\right)  \otimes%
%TCIMACRO{\tbigwedge }%
%BeginExpansion
{\textstyle\bigwedge}
%EndExpansion
\left(  \mathcal{H}_{\mathcal{F}}^{d}\right)  }\right)
\end{equation}
and thus use the same symbols $\bullet$ and $\boxtimes$ for the products in
$\mathcal{B}_{00}\left(  \mathcal{H}_{\mathcal{F}}\right)  $ and $%
%TCIMACRO{\tbigwedge }%
%BeginExpansion
{\textstyle\bigwedge}
%EndExpansion
\left(  \mathcal{H}_{\mathcal{F}}\right)  \otimes%
%TCIMACRO{\tbigwedge }%
%BeginExpansion
{\textstyle\bigwedge}
%EndExpansion
\left(  \mathcal{H}_{\mathcal{F}}^{d}\right)  .$
\end{notation}

We show later that the $\boxtimes_{W}$product of projector onto $\mathcal{H}%
^{1}$ is given by
\begin{gather}
\left(  \sum_{1\leq j_{1}\leq n}\left\vert x_{j_{1}}\right\rangle \left\langle
x_{j_{1}}\right\vert \right)  \boxtimes_{W}\left(  \sum_{1\leq k_{1}\leq
n}\left\vert x_{k_{1}}\right\rangle \left\langle x_{k_{1}}\right\vert \right)
=\left(  \sum_{1\leq j_{1},k_{1}\leq n}\left\vert x_{j_{1}}\wedge_{W}x_{k_{1}%
}\right\rangle \left\langle x_{j_{1}}\wedge_{W}x_{k_{1}}\right\vert \right)
=2\left(  \sum_{1\leq j_{1}<k_{1}\leq n}\left\vert x_{j_{1}}\wedge_{W}%
x_{k_{1}}\right\rangle \left\langle x_{j_{1}}\wedge_{W}x_{k_{1}}\right\vert
\right) \\
=2\frac{1}{2}\left(  \sum_{1\leq j_{1}<k_{1}\leq n}\left\vert x_{j_{1}%
}x_{k_{1}}\right\rangle \left\langle x_{j_{1}}x_{k_{1}}\right\vert \right)
=\sum_{1\leq j_{1}<k_{1}\leq n}\left\vert x_{j_{1}}x_{k_{1}}\right\rangle
\left\langle x_{j_{1}}x_{k_{1}}\right\vert
\end{gather}
i.e.
\begin{equation}
P_{\mathcal{H}^{1}}\boxtimes_{W}P_{\mathcal{H}^{1}}=P_{\mathcal{H}^{2}%
}=P_{\mathcal{H}^{1}}^{2\boxtimes_{W}}.
\end{equation}
On the other hand using the Greub product $\boxtimes_{G}$ for $n=4$ one
obtains
\begin{align}
&  \left(  \left\vert x_{1}\right\rangle \left\langle x_{1}\right\vert
+\left\vert x_{2}\right\rangle \left\langle x_{2}\right\vert +\left\vert
x_{3}\right\rangle \left\langle x_{3}\right\vert +\left\vert x_{4}%
\right\rangle \left\langle x_{4}\right\vert \right)  \boxtimes_{G}\left(
\left\vert x_{1}\right\rangle \left\langle x_{1}\right\vert +\left\vert
x_{2}\right\rangle \left\langle x_{2}\right\vert +\left\vert x_{3}%
\right\rangle \left\langle x_{3}\right\vert +\left\vert x_{4}\right\rangle
\left\langle x_{4}\right\vert \right) \\
&  =\left\vert x_{1}\right\rangle \left\langle x_{1}\right\vert \boxtimes
_{G}\left(  \left\vert x_{2}\right\rangle \left\langle x_{2}\right\vert
+\left\vert x_{3}\right\rangle \left\langle x_{3}\right\vert +\left\vert
x_{4}\right\rangle \left\langle x_{4}\right\vert \right)  +\left\vert
x_{2}\right\rangle \left\langle x_{2}\right\vert \boxtimes_{G}\left(
\left\vert x_{1}\right\rangle \left\langle x_{1}\right\vert +\left\vert
x_{3}\right\rangle \left\langle x_{3}\right\vert +\left\vert x_{4}%
\right\rangle \left\langle x_{4}\right\vert \right) \\
&  +\left\vert x_{3}\right\rangle \left\langle x_{3}\right\vert \boxtimes
_{G}\left(  \left\vert x_{2}\right\rangle \left\langle x_{2}\right\vert
+\left\vert x_{1}\right\rangle \left\langle x_{1}\right\vert +\left\vert
x_{4}\right\rangle \left\langle x_{4}\right\vert \right)  +\left\vert
x_{4}\right\rangle \left\langle x_{4}\right\vert \boxtimes_{G}\left(
\left\vert x_{1}\right\rangle \left\langle x_{1}\right\vert +\left\vert
x_{2}\right\rangle \left\langle x_{2}\right\vert +\left\vert x_{3}%
\right\rangle \left\langle x_{3}\right\vert \right) \\
&  =2\left(  \left\vert x_{1}\wedge_{G}x_{2}\right\rangle \left\langle
x_{1}\wedge_{G}x_{2}\right\vert +\left\vert x_{1}\wedge_{G}x_{3}\right\rangle
\left\langle x_{1}\wedge_{G}x_{3}\right\vert +\left\vert x_{1}\wedge_{G}%
x_{4}\right\rangle \left\langle x_{1}\wedge_{G}x_{4}\right\vert \right) \\
&  +2\left(  \left\vert x_{2}\wedge_{G}x_{3}\right\rangle \left\langle
x_{2}\wedge_{G}x_{3}\right\vert +\left\vert x_{2}\wedge_{G}x_{4}\right\rangle
\left\langle x_{2}\wedge_{G}x_{4}\right\vert \right)  +2\left\vert x_{3}%
\wedge_{G}x_{4}\right\rangle \left\langle x_{3}\wedge_{G}x_{4}\right\vert \\
&  =4\left(  \left\vert x_{1}x_{2}\right\rangle \left\langle x_{1}%
x_{2}\right\vert +\left\vert x_{1}x_{3}\right\rangle \left\langle x_{1}%
x_{3}\right\vert +\left\vert x_{1}x_{4}\right\rangle \left\langle x_{1}%
x_{4}\right\vert +\left\vert x_{2}x_{3}\right\rangle \left\langle x_{2}%
x_{3}\right\vert +\left\vert x_{3}x_{4}\right\rangle \left\langle x_{3}%
x_{4}\right\vert \right)  ,
\end{align}
i.e.
\begin{equation}
P_{\mathcal{H}^{1}}\boxtimes_{G}P_{\mathcal{H}^{1}}=2P_{\mathcal{H}^{1}%
}^{2\boxtimes_{G}}=\binom{2}{1}P_{\mathcal{H}^{1}}^{2\boxtimes_{G}%
}=4P_{\mathcal{H}^{2}}=\left(  2!\right)  ^{2}P_{\mathcal{H}^{2}}.
\end{equation}


\begin{example}
The $\boxtimes_{W}$product of projectors $P_{\mathcal{H}^{1}}$ gives
$P_{\mathcal{H}^{2}}$%
\begin{align}
&  \left(  \left\vert x_{1}\right\rangle \left\langle x_{1}\right\vert
+\left\vert x_{2}\right\rangle \left\langle x_{2}\right\vert +\left\vert
x_{3}\right\rangle \left\langle x_{3}\right\vert +\left\vert x_{4}%
\right\rangle \left\langle x_{4}\right\vert \right)  \boxtimes_{W}\left(
\begin{array}
[c]{c}%
\left\vert x_{1}x_{2}\right\rangle \left\langle x_{1}x_{2}\right\vert
+\left\vert x_{1}x_{3}\right\rangle \left\langle x_{1}x_{3}\right\vert
+\left\vert x_{1}x_{4}\right\rangle \left\langle x_{1}x_{4}\right\vert \\
+\left\vert x_{2}x_{3}\right\rangle \left\langle x_{2}x_{3}\right\vert
+\left\vert x_{2}x_{4}\right\rangle \left\langle x_{2}x_{4}\right\vert
+\left\vert x_{3}x_{4}\right\rangle \left\langle x_{3}x_{4}\right\vert
\end{array}
\right) \\
&  =\left\vert x_{1}\right\rangle \left\langle x_{1}\right\vert \boxtimes
_{W}\left(  \left\vert x_{2}x_{3}\right\rangle \left\langle x_{2}%
x_{3}\right\vert +\left\vert x_{2}x_{4}\right\rangle \left\langle x_{2}%
x_{4}\right\vert +\left\vert x_{3}x_{4}\right\rangle \left\langle x_{3}%
x_{4}\right\vert \right)  +\left\vert x_{2}\right\rangle \left\langle
x_{2}\right\vert \boxtimes_{W}\left(  \left\vert x_{1}x_{3}\right\rangle
\left\langle x_{1}x_{3}\right\vert +\left\vert x_{1}x_{4}\right\rangle
\left\langle x_{1}x_{4}\right\vert +\left\vert x_{3}x_{4}\right\rangle
\left\langle x_{3}x_{4}\right\vert \right) \\
&  +\left\vert x_{3}\right\rangle \left\langle x_{3}\right\vert \boxtimes
_{W}\left(  \left\vert x_{1}x_{2}\right\rangle \left\langle x_{1}%
x_{2}\right\vert +\left\vert x_{2}x_{4}\right\rangle \left\langle x_{2}%
x_{4}\right\vert +\left\vert x_{2}x_{4}\right\rangle \left\langle x_{2}%
x_{4}\right\vert \right)  +\left\vert x_{4}\right\rangle \left\langle
x_{4}\right\vert \boxtimes_{W}\left(  \left\vert x_{1}x_{2}\right\rangle
\left\langle x_{1}x_{2}\right\vert +\left\vert x_{1}x_{3}\right\rangle
\left\langle x_{1}x_{3}\right\vert +\left\vert x_{2}x_{3}\right\rangle
\left\langle x_{2}x_{3}\right\vert \right) \\
&  =2\left\vert x_{1}\right\rangle \left\langle x_{1}\right\vert \boxtimes
_{W}\left(  \left\vert x_{2}\wedge_{W}x_{3}\right\rangle \left\langle
x_{2}\wedge_{W}x_{3}\right\vert +\left\vert x_{2}\wedge_{W}x_{4}\right\rangle
\left\langle x_{1}\wedge_{W}x_{2}\right\vert +\left\vert x_{3}\wedge_{W}%
x_{4}\right\rangle \left\langle x_{3}\wedge_{W}x_{4}\right\vert \right) \\
&  +2\left\vert x_{2}\right\rangle \left\langle x_{2}\right\vert \boxtimes
_{W}\left(  \left\vert x_{1}\wedge_{W}x_{3}\right\rangle \left\langle
x_{1}\wedge_{W}x_{3}\right\vert +\left\vert x_{1}\wedge_{W}x_{4}\right\rangle
\left\langle x_{1}\wedge_{W}x_{4}\right\vert +\left\vert x_{3}\wedge_{W}%
x_{4}\right\rangle \left\langle x_{3}\wedge_{W}x_{4}\right\vert \right) \\
&  +2\left\vert x_{3}\right\rangle \left\langle x_{3}\right\vert \boxtimes
_{W}\left(  \left\vert x_{1}\wedge_{W}x_{2}\right\rangle \left\langle
x_{1}\wedge_{W}x_{2}\right\vert +\left\vert x_{2}\wedge_{W}x_{4}\right\rangle
\left\langle x_{2}\wedge_{W}x_{4}\right\vert +\left\vert x_{1}\wedge_{W}%
x_{4}\right\rangle \left\langle x_{1}\wedge_{W}x_{4}\right\vert \right) \\
&  +2\left\vert x_{4}\right\rangle \left\langle x_{4}\right\vert \boxtimes
_{W}\left(  \left\vert x_{1}\wedge_{W}x_{2}\right\rangle \left\langle
x_{1}\wedge_{W}x_{2}\right\vert +\left\vert x_{1}\wedge_{W}x_{3}\right\rangle
\left\langle x_{1}\wedge_{W}x_{3}\right\vert +\left\vert x_{2}\wedge_{W}%
x_{3}\right\rangle \left\langle x_{2}\wedge_{W}x_{3}\right\vert \right) \\
&  =\left(  2\times3\right)  \left(
\begin{array}
[c]{c}%
\left\vert x_{1}\wedge_{W}x_{2}\wedge_{W}x_{3}\right\rangle \left\langle
x_{1}\wedge_{W}x_{2}\wedge_{W}x_{3}\right\vert +\left\vert x_{1}\wedge
_{W}x_{2}\wedge_{W}x_{4}\right\rangle \left\langle x_{1}\wedge_{W}x_{2}%
\wedge_{W}x_{4}\right\vert \\
+\left\vert x_{1}\wedge_{W}x_{3}\wedge_{W}x_{4}\right\rangle \left\langle
x_{1}\wedge_{W}x_{3}\wedge_{W}x_{4}\right\vert +\left\vert x_{2}\wedge
_{W}x_{3}\wedge_{W}x_{4}\right\rangle \left\langle x_{2}\wedge_{W}x_{3}%
\wedge_{W}x_{4}\right\vert
\end{array}
\right) \\
&  =\left\vert x_{1}x_{2}x_{3}\right\rangle \left\langle x_{1}x_{2}%
x_{3}\right\vert +\left\vert x_{1}x_{2}x_{4}\right\rangle \left\langle
x_{1}x_{2}x_{4}\right\vert +\left\vert x_{1}x_{3}x_{4}\right\rangle
\left\langle x_{1}x_{3}x_{4}\right\vert +\left\vert x_{2}x_{3}x_{4}%
\right\rangle \left\langle x_{2}x_{3}x_{4}\right\vert
\end{align}
hence
\begin{equation}
P_{\mathcal{H}^{1}}\boxtimes_{W}P_{\mathcal{H}^{2}}=P_{\mathcal{H}^{3}}.
\end{equation}
The $\boxtimes_{G}$product of projectors $P_{\mathcal{H}^{1}}$ gives
$P_{\mathcal{H}^{2}}$%
\begin{align}
&  \left(  \left\vert x_{1}\right\rangle \left\langle x_{1}\right\vert
+\left\vert x_{2}\right\rangle \left\langle x_{2}\right\vert +\left\vert
x_{3}\right\rangle \left\langle x_{3}\right\vert +\left\vert x_{4}%
\right\rangle \left\langle x_{4}\right\vert \right)  \boxtimes_{G}\left(
\begin{array}
[c]{c}%
\left\vert x_{1}x_{2}\right\rangle \left\langle x_{1}x_{2}\right\vert
+\left\vert x_{1}x_{3}\right\rangle \left\langle x_{1}x_{3}\right\vert
+\left\vert x_{1}x_{4}\right\rangle \left\langle x_{1}x_{4}\right\vert \\
+\left\vert x_{2}x_{3}\right\rangle \left\langle x_{2}x_{3}\right\vert
+\left\vert x_{2}x_{4}\right\rangle \left\langle x_{2}x_{4}\right\vert
+\left\vert x_{3}x_{4}\right\rangle \left\langle x_{3}x_{4}\right\vert
\end{array}
\right) \\
&  =\left\vert x_{1}\right\rangle \left\langle x_{1}\right\vert \boxtimes
_{G}\left(  \left\vert x_{2}x_{3}\right\rangle \left\langle x_{2}%
x_{3}\right\vert +\left\vert x_{2}x_{4}\right\rangle \left\langle x_{2}%
x_{4}\right\vert +\left\vert x_{3}x_{4}\right\rangle \left\langle x_{3}%
x_{4}\right\vert \right)  +\left\vert x_{2}\right\rangle \left\langle
x_{2}\right\vert \boxtimes_{G}\left(  \left\vert x_{1}x_{3}\right\rangle
\left\langle x_{1}x_{3}\right\vert +\left\vert x_{1}x_{4}\right\rangle
\left\langle x_{1}x_{4}\right\vert +\left\vert x_{3}x_{4}\right\rangle
\left\langle x_{3}x_{4}\right\vert \right) \\
&  +\left\vert x_{3}\right\rangle \left\langle x_{3}\right\vert \boxtimes
_{G}\left(  \left\vert x_{1}x_{2}\right\rangle \left\langle x_{1}%
x_{2}\right\vert +\left\vert x_{2}x_{4}\right\rangle \left\langle x_{2}%
x_{4}\right\vert +\left\vert x_{1}x_{4}\right\rangle \left\langle x_{1}%
x_{4}\right\vert \right)  +\left\vert x_{4}\right\rangle \left\langle
x_{4}\right\vert \boxtimes_{G}\left(  \left\vert x_{1}x_{2}\right\rangle
\left\langle x_{1}x_{2}\right\vert +\left\vert x_{1}x_{3}\right\rangle
\left\langle x_{1}x_{3}\right\vert +\left\vert x_{2}x_{3}\right\rangle
\left\langle x_{2}x_{3}\right\vert \right) \\
&  =\frac{1}{4}\left\vert x_{1}\right\rangle \left\langle x_{1}\right\vert
\boxtimes_{G}\left(  \left\vert x_{2}\wedge_{G}x_{3}\right\rangle \left\langle
x_{2}\wedge_{G}x_{3}\right\vert +\left\vert x_{2}\wedge_{G}x_{4}\right\rangle
\left\langle x_{2}\wedge_{G}x_{4}\right\vert +\left\vert x_{3}\wedge_{G}%
x_{4}\right\rangle \left\langle x_{3}\wedge_{G}x_{4}\right\vert \right) \\
&  +\frac{1}{4}\left\vert x_{2}\right\rangle \left\langle x_{2}\right\vert
\boxtimes_{G}\left(  \left\vert x_{1}\wedge_{G}x_{3}\right\rangle \left\langle
x_{1}\wedge_{G}x_{3}\right\vert +\left\vert x_{1}\wedge_{G}x_{4}\right\rangle
\left\langle x_{1}\wedge_{G}x_{4}\right\vert +\left\vert x_{3}\wedge_{G}%
x_{4}\right\rangle \left\langle x_{3}\wedge_{G}x_{4}\right\vert \right) \\
&  +\frac{1}{4}\left\vert x_{3}\right\rangle \left\langle x_{3}\right\vert
\boxtimes_{G}\left(  \left\vert x_{1}\wedge_{G}x_{2}\right\rangle \left\langle
x_{1}\wedge_{G}x_{2}\right\vert +\left\vert x_{2}\wedge_{G}x_{4}\right\rangle
\left\langle x_{2}\wedge_{G}x_{4}\right\vert +\left\vert x_{1}\wedge_{G}%
x_{4}\right\rangle \left\langle x_{1}\wedge_{G}x_{4}\right\vert \right) \\
&  +\frac{1}{4}\left\vert x_{4}\right\rangle \left\langle x_{4}\right\vert
\boxtimes_{G}\left(  \left\vert x_{1}\wedge_{G}x_{2}\right\rangle \left\langle
x_{1}\wedge_{G}x_{2}\right\vert +\left\vert x_{1}\wedge_{G}x_{3}\right\rangle
\left\langle x_{1}\wedge_{G}x_{3}\right\vert +\left\vert x_{2}\wedge_{G}%
x_{3}\right\rangle \left\langle x_{2}\wedge_{G}x_{3}\right\vert \right) \\
&  =\frac{1}{4}\left(  \left\vert x_{1}\wedge_{G}x_{2}\wedge_{G}%
x_{3}\right\rangle \left\langle x_{1}\wedge_{G}x_{2}\wedge_{G}x_{3}\right\vert
+\left\vert x_{1}\wedge_{G}x_{2}\wedge_{G}x_{4}\right\rangle \left\langle
x_{1}\wedge_{G}x_{2}\wedge_{G}x_{4}\right\vert +\left\vert x_{1}\wedge
_{G}x_{3}\wedge_{G}x_{4}\right\rangle \left\langle x_{1}\wedge_{G}x_{3}%
\wedge_{G}x_{4}\right\vert \right. \\
&  +\left\vert x_{2}\wedge_{G}x_{1}\wedge_{G}x_{3}\right\rangle \left\langle
x_{2}\wedge_{G}x_{1}\wedge_{G}x_{3}\right\vert +\left\vert x_{2}\wedge
_{G}x_{1}\wedge_{G}x_{4}\right\rangle \left\langle x_{2}\wedge_{G}x_{1}%
\wedge_{G}x_{4}\right\vert +\left\vert x_{2}\wedge_{G}x_{3}\wedge_{G}%
x_{4}\right\rangle \left\langle x_{2}\wedge_{G}x_{3}\wedge_{G}x_{4}\right\vert
\\
&  +\left\vert x_{3}\wedge_{G}x_{1}\wedge_{G}x_{2}\right\rangle \left\langle
x_{3}\wedge_{G}x_{1}\wedge_{G}x_{2}\right\vert +\left\vert x_{3}\wedge
_{G}x_{2}\wedge_{G}x_{4}\right\rangle \left\langle x_{3}\wedge_{G}x_{2}%
\wedge_{G}x_{4}\right\vert +\left\vert x_{3}\wedge_{G}x_{1}\wedge_{G}%
x_{4}\right\rangle \left\langle x_{3}\wedge_{G}x_{1}\wedge_{G}x_{4}\right\vert
\\
&  +\left.  \left\vert x_{4}\wedge_{G}x_{1}\wedge_{G}x_{2}\right\rangle
\left\langle x_{4}\wedge_{G}x_{1}\wedge_{G}x_{2}\right\vert +\left\vert
x_{4}\wedge_{G}x_{1}\wedge_{G}x_{3}\right\rangle \left\langle x_{4}\wedge
_{G}x_{1}\wedge_{G}x_{3}\right\vert +\left\vert x_{4}\wedge_{G}x_{2}\wedge
_{G}x_{3}\right\rangle \left\langle x_{4}\wedge_{G}x_{2}\wedge_{G}%
x_{3}\right\vert \right)
\end{align}%
\begin{align}
&  =\frac{3}{4}\left(  \left\vert x_{1}\wedge_{G}x_{2}\wedge_{G}%
x_{3}\right\rangle \left\langle x_{1}\wedge_{G}x_{2}\wedge_{G}x_{3}\right\vert
+\left\vert x_{1}\wedge_{G}x_{2}\wedge_{G}x_{4}\right\rangle \left\langle
x_{1}\wedge_{G}x_{2}\wedge_{G}x_{4}\right\vert +\left\vert x_{1}\wedge
_{G}x_{3}\wedge_{G}x_{4}\right\rangle \left\langle x_{1}\wedge_{G}x_{3}%
\wedge_{G}x_{4}\right\vert \right. \\
&  +\left.  \left\vert x_{2}\wedge_{G}x_{3}\wedge_{G}x_{4}\right\rangle
\left\langle x_{2}\wedge_{G}x_{3}\wedge_{G}x_{4}\right\vert \right) \\
&  \frac{3}{4}\left(  3!\right)  ^{2}\left(  \left\vert x_{1}x_{2}%
x_{3}\right\rangle \left\langle x_{1}x_{2}x_{3}\right\vert +\left\vert
x_{1}x_{2}x_{4}\right\rangle \left\langle x_{1}x_{2}x_{4}\right\vert
+\left\vert x_{1}x_{3}x_{4}\right\rangle \left\langle x_{1}x_{3}%
x_{4}\right\vert \right. \\
&  +\left.  \left\vert x_{2}x_{3}x_{4}\right\rangle \left\langle x_{2}%
x_{3}x_{4}\right\vert \right)
\end{align}
i.e
\begin{equation}
P_{\mathcal{H}^{1}}\boxtimes_{G}P_{\mathcal{H}^{2}}=\frac{3}{4}\left(
3!\right)  ^{2}P_{\mathcal{H}^{3}}=\binom{3}{1}P_{\mathcal{H}^{1}}%
^{3\boxtimes_{G}}%
\end{equation}
$\lceil$%
\begin{align}
P_{\mathcal{H}^{1}}\boxtimes_{G}P_{\mathcal{H}^{2}}  &  =\frac{3}{4}\left(
3!\right)  ^{2}P_{\mathcal{H}^{3}}.\\
P_{\mathcal{H}^{1}}\boxtimes_{G}\frac{1}{4}P_{\mathcal{H}^{1}}^{2\boxtimes
_{G}}  &  =\frac{3}{4}\left(  3!\right)  ^{2}P_{\mathcal{H}^{3}}\\
P_{\mathcal{H}^{1}}\boxtimes_{G}P_{\mathcal{H}^{1}}^{2\boxtimes_{G}}  &
=P_{\mathcal{H}^{1}}^{3\boxtimes_{G}}=3\left(  3!\right)  ^{2}P_{\mathcal{H}%
^{3}}\\
P_{\mathcal{H}^{1}}^{3\boxtimes_{G}}  &  =\left(  3!\right)  ^{2}%
P_{\mathcal{H}^{3}}\\
P_{\mathcal{H}^{3}}  &  =\left(  3!\right)  ^{-2}P_{\mathcal{H}^{1}%
}^{3\boxtimes_{G}}\\
P_{\mathcal{H}^{1}}\boxtimes_{G}P_{\mathcal{H}^{1}}^{2\boxtimes_{G}}  &
=P_{\mathcal{H}^{1}}^{3\boxtimes_{G}}=3\left(  3!\right)  ^{2}\left(
3!\right)  ^{-2}P_{\mathcal{H}^{1}}^{3\boxtimes_{G}}=3P_{\mathcal{H}^{1}%
}^{3\boxtimes_{G}}=\binom{3}{1}P_{\mathcal{H}^{1}}^{3\boxtimes_{G}}\\
P_{\mathcal{H}^{1}}^{2\boxtimes_{G}}  &  =\left(  2!\right)  ^{2}%
P_{\mathcal{H}^{2}}%
\end{align}
$\rfloor$
\end{example}

\subsection{Duality}

Let $\mathcal{H},\mathcal{H}^{d}$ be dual spaces and let $T$ be an injective
inclusion of $\mathcal{H}$ into $\mathcal{H}^{d}$
\begin{equation}
T:\mathcal{H}\rightarrow\mathcal{H}^{d}%
\end{equation}
then if
\begin{equation}
A^{d}:\mathcal{H}^{d}\rightarrow\mathcal{H}^{d}%
\end{equation}
then
\begin{equation}
T^{-1}A^{d}T:\mathcal{H}\rightarrow\mathcal{H}.
\end{equation}


\begin{remark}
In general there is \emph{no }map between $\mathcal{H}$ and $\mathcal{H}^{d}.$
Of course there is always an injection of $\mathcal{H}$ into $\mathcal{H}%
^{dd}.$
\end{remark}

In the particular case of a Hilbert space $\mathcal{H},\mathcal{H}%
\overset{anti}{\mathcal{\cong}}\mathcal{H}^{d}$ (antilinear isometric
isomorphism) and $T$ can be identified as the adjoint operation "$\dagger$"
i.e.%
\begin{equation}
T\left(  \left\vert \Psi\right\rangle \right)  =\left\langle \Psi\right\vert
=\left\vert \Psi\right\rangle ^{\dagger}\equiv\left\vert \Psi\right\rangle
^{d}.
\end{equation}
The adjoint operation acts in both $\mathcal{H}$ and $\mathcal{H}^{d}$ i.e.%
\begin{align}
\dagger &  :\mathcal{H}\rightarrow\mathcal{H}^{d}\\
\dagger &  :\mathcal{H}^{d}\rightarrow\mathcal{H}%
\end{align}
and
\begin{equation}
\dagger^{2}=I_{\mathcal{H}}=I_{\mathcal{H}^{d}}.
\end{equation}


\begin{remark}
$\mathcal{H}$ is a reflexive Banach space.
\end{remark}

The adjoint map extends to $\mathcal{B}\left(  \mathcal{H}\right)  $ and
gives
\begin{equation}
TAT^{-1}=A^{\#}.
\end{equation}
The action of $\dagger$ in $\mathcal{B}\left(  \mathcal{H}\right)  $ can be
expressed as
\begin{equation}
A=\sum_{1\leq j,k\leq n}A_{jk}\left\vert \Psi_{k}\right\rangle \left\langle
\Psi_{j}\right\vert \in\mathcal{B}\left(  \mathcal{H}\right)  =\mathcal{B}%
\left(  \mathcal{H}^{d}\right)
\end{equation}
and
\begin{equation}
A^{\dagger}=\sum_{1\leq j,k\leq n}\overline{A_{jk}}\left\vert \Psi
_{j}\right\rangle \left\langle \Psi_{k}\right\vert \in\mathcal{B}\left(
\mathcal{H}^{d}\right)  =\mathcal{B}\left(  \mathcal{H}\right)  ,
\end{equation}
as
\begin{equation}
\left(  \left\vert \Psi_{k}\right\rangle \left\langle \Psi_{j}\right\vert
\right)  ^{\dagger}=\left(  \left\langle \Psi_{j}\right\vert \right)
^{^{\dagger}}\left(  \left\vert \Psi_{k}\right\rangle \right)  ^{^{\dagger}%
}=\left\vert \Psi_{j}\right\rangle \left\langle \Psi_{k}\right\vert .
\end{equation}
Hence
\begin{equation}
T\left(  A\left\vert \Psi\right\rangle \right)  =TAT^{-1}T\left\vert
\Psi\right\rangle =A^{\#}T\left(  \left\vert \Psi\right\rangle \right)
=\left\langle \Psi\right\vert A^{\dagger}=\left(  A\left\vert \Psi
\right\rangle \right)  ^{\dagger}.
\end{equation}


\begin{remark}
\emph{Note that the spaces of operators on }$\mathcal{H}$ and $\mathcal{H}%
^{d}$ are antilinear isomorphic i.e. \emph{ }
\begin{equation}
\mathcal{B}\left(  \mathcal{H}\right)  \overset{anti}{\mathcal{\cong}%
}\mathcal{B}\left(  \mathcal{H}^{d}\right)  ,
\end{equation}
but they are \emph{equal as }$c^{\ast}$-algebras and are self dual.
\end{remark}

\begin{example}
Consider the maps
\end{example}

%

\begin{align}
A  &  :\mathcal{H}\rightarrow\mathcal{H}\\
A^{d}  &  :\mathcal{H}^{d}\rightarrow\mathcal{H}^{d}%
\end{align}
and the isomorphisms
\begin{align}
\dagger &  :\mathcal{H}\rightarrow\mathcal{H}^{d}\\
\dagger &  :\mathcal{H}^{d}\rightarrow\mathcal{H}%
\end{align}
then
\begin{align}
A\left\vert \Psi\right\rangle  &  =\left\vert \Phi\right\rangle \\
\left(  A\left\vert \Psi\right\rangle \right)  ^{\dagger}  &  =\left(
\left\vert \Phi\right\rangle \right)  ^{\dagger}\\
\left\langle \Psi\right\vert A^{\dagger}  &  =\left\langle \Phi\right\vert
\end{align}
and%
\begin{align}
\left\langle \Psi\right\vert B  &  =\left\langle \Phi\right\vert \\
\left(  \left\langle \Psi\right\vert B\right)  ^{\dagger}  &  =\left(
\left\langle \Phi\right\vert \right)  ^{\dagger}\\
B^{\dagger}\left\vert \Psi\right\rangle  &  =\left\vert \Phi\right\rangle .
\end{align}


\begin{remark}%
\begin{equation}
\left(  \left\vert y_{1}\right\rangle \wedge\left\vert y_{2}\right\rangle
\right)  ^{\dagger}\equiv\left(  \left\vert y_{1}\wedge y_{2}\right\rangle
\right)  ^{\dagger}=\left\langle y_{1}\wedge y_{2}\right\vert =\left\langle
y_{1}\right\vert \wedge\left\langle y_{2}\right\vert .
\end{equation}

\end{remark}

\subsection{Mixed Exterior Algebra}

\begin{remark}
In the following we denote $\mathcal{H}$ as $\mathcal{H}^{1}$ and the exterior
algebra generated by $\mathcal{H}^{1}$ as $\mathcal{H}_{\mathcal{E}}=%
%TCIMACRO{\tbigwedge }%
%BeginExpansion
{\textstyle\bigwedge}
%EndExpansion
\left(  \mathcal{H}^{1}\right)  ,$ which is also a vector space.
\end{remark}

The mixed exterior algebra, $\mathcal{E},$ over $\left\{  \mathcal{H}%
^{d},\mathcal{H}\right\}  $ is defined as
\begin{equation}
\mathcal{E}=%
%TCIMACRO{\tbigwedge }%
%BeginExpansion
{\textstyle\bigwedge}
%EndExpansion
\left(  \mathcal{H}^{1},\mathcal{H}^{1d}\right)  =%
%TCIMACRO{\tbigwedge }%
%BeginExpansion
{\textstyle\bigwedge}
%EndExpansion
\left(  \mathcal{H}^{1}\right)  \otimes%
%TCIMACRO{\tbigwedge }%
%BeginExpansion
{\textstyle\bigwedge}
%EndExpansion
\left(  \mathcal{H}^{1d}\right)  =\mathcal{H}_{\mathcal{E}}\otimes
\mathcal{H}_{\mathcal{E}}^{d},
\end{equation}
The duality is expressed by
\begin{equation}
\left\langle ,\right\rangle :\mathcal{H}_{\mathcal{E}}\times\mathcal{H}%
_{\mathcal{E}}^{d}\rightarrow\mathbb{C}.
\end{equation}


\begin{remark}
The mixed exterior algebra $%
%TCIMACRO{\tbigwedge }%
%BeginExpansion
{\textstyle\bigwedge}
%EndExpansion
\left(  \mathcal{H}^{1}\right)  \otimes%
%TCIMACRO{\tbigwedge }%
%BeginExpansion
{\textstyle\bigwedge}
%EndExpansion
\left(  \mathcal{H}^{1d}\right)  $ is isomorphic to the bounded operators
$\mathcal{B}\left(  \mathcal{H}_{\mathcal{F}}\right)  $ on $\mathcal{H}%
_{\mathcal{E}}=\mathcal{H}_{\mathcal{F}}=%
%TCIMACRO{\tbigwedge }%
%BeginExpansion
{\textstyle\bigwedge}
%EndExpansion
\left(  \mathcal{H}^{1}\right)  ,$ when one takes operator norm limits of
nuclear operators i.e.
\begin{equation}
\mathcal{F}=\mathcal{B}\left(  \mathcal{H}_{\mathcal{F}}\right)  \cong%
%TCIMACRO{\tbigwedge }%
%BeginExpansion
{\textstyle\bigwedge}
%EndExpansion
\left(  \mathcal{H}^{1},\mathcal{H}^{1d}\right)  .
\end{equation}
The isomorphism is denoted by $\Xi_{\mathcal{H}_{\mathcal{F}}\mathcal{H}%
_{\mathcal{F}}^{d}}$
\begin{equation}
\mathcal{F}=\mathcal{B}\left(  \mathcal{H}_{\mathcal{F}}\right)
=\Xi_{\mathcal{H}_{\mathcal{F}}\mathcal{H}_{\mathcal{F}}^{d}}\left(
%TCIMACRO{\tbigwedge }%
%BeginExpansion
{\textstyle\bigwedge}
%EndExpansion
\left(  \mathcal{H}^{1},\mathcal{H}^{1d}\right)  \right)  =\Xi_{\mathcal{H}%
_{\mathcal{F}}\mathcal{H}_{\mathcal{F}}^{d}}\left(
%TCIMACRO{\tbigwedge }%
%BeginExpansion
{\textstyle\bigwedge}
%EndExpansion
\left(  \mathcal{H}^{1}\right)  \otimes%
%TCIMACRO{\tbigwedge }%
%BeginExpansion
{\textstyle\bigwedge}
%EndExpansion
\left(  \mathcal{H}^{1d}\right)  \right)  =\Xi_{\mathcal{H}_{\mathcal{F}%
}\mathcal{H}_{\mathcal{F}}^{d}}\left(  \mathcal{H}_{\mathcal{F}}%
\otimes\mathcal{H}_{\mathcal{F}}^{d}\right)  ,
\end{equation}
where
\begin{equation}
\mathcal{H}_{\mathcal{F}}=%
%TCIMACRO{\tbigwedge }%
%BeginExpansion
{\textstyle\bigwedge}
%EndExpansion
\left(  \mathcal{H}^{1}\right)
\end{equation}
and
\begin{equation}
\mathcal{H}_{\mathcal{F}}^{d}=%
%TCIMACRO{\tbigwedge }%
%BeginExpansion
{\textstyle\bigwedge}
%EndExpansion
\left(  \mathcal{H}^{1d}\right)  .
\end{equation}

\end{remark}

\begin{remark}
Greub considers $%
%TCIMACRO{\tbigwedge }%
%BeginExpansion
{\textstyle\bigwedge}
%EndExpansion
\left(  \mathcal{H}^{1d},\mathcal{H}^{1}\right)  $ as the mixed exterior
algebra, however $%
%TCIMACRO{\tbigwedge }%
%BeginExpansion
{\textstyle\bigwedge}
%EndExpansion
\left(  \mathcal{H}^{1d},\mathcal{H}^{1}\right)  \cong$ $%
%TCIMACRO{\tbigwedge }%
%BeginExpansion
{\textstyle\bigwedge}
%EndExpansion
\left(  \mathcal{H}^{1},\mathcal{H}^{1d}\right)  $
\end{remark}

We represent this by the ket-bra formalism
\begin{equation}
\left\vert \Psi\right\rangle \left\langle \Phi\right\vert \in\mathcal{\Xi
}_{\mathcal{H}_{\mathcal{F}}\mathcal{H}_{\mathcal{F}}^{d}}\left(
%TCIMACRO{\tbigwedge }%
%BeginExpansion
{\textstyle\bigwedge}
%EndExpansion
\left(  \mathcal{H}^{1},\mathcal{H}^{1d}\right)  \right)  =\mathcal{\Xi
}_{\mathcal{H}_{\mathcal{F}}\mathcal{H}_{\mathcal{F}}^{d}}\left(
%TCIMACRO{\tbigwedge }%
%BeginExpansion
{\textstyle\bigwedge}
%EndExpansion
\left(  \mathcal{H}^{1}\right)  \otimes%
%TCIMACRO{\tbigwedge }%
%BeginExpansion
{\textstyle\bigwedge}
%EndExpansion
\left(  \mathcal{H}^{1d}\right)  \right)  ;\,\left\vert \Psi\right\rangle
\in\mathcal{H}_{\mathcal{F}},\left\langle \Phi\right\vert \in\mathcal{H}%
_{\mathcal{F}}^{d}.
\end{equation}
The map between $\mathcal{H}$ and $\mathcal{H}^{d}$ is given by
\begin{equation}
\dagger:\mathcal{H}_{\mathcal{F}}\rightarrow\mathcal{H}_{\mathcal{F}}^{d}%
\end{equation}
i.e.
\begin{equation}
\left(  \left\vert \Psi\right\rangle \right)  ^{\dagger}\equiv\left\vert
\Psi\right\rangle ^{d}\rightarrow\left\langle \Psi\right\vert .
\end{equation}
This can be made rigorous by \emph{rigging}. Let
\begin{equation}
v^{d}=v^{\dagger}%
\end{equation}
then
\begin{equation}
\left\langle v^{d},w\right\rangle =\left\langle v^{\dagger},w\right\rangle
=\left\langle \left.  v\right\vert w\right\rangle ,
\end{equation}
where $\left\langle \left.  v\right\vert w\right\rangle $ is the inner product
in $\mathcal{H}^{1}$ if it is a Hilbert space.

\begin{definition}
The isomorphism $\mathcal{\Xi}_{\mathcal{H}_{\mathcal{F}}\mathcal{H}%
_{\mathcal{F}}^{d}}$ between the mixed exterior algebra generated
$\mathcal{H}_{\mathcal{F}}$ by and the space of bounded operators acting in
$\mathcal{B}\left(  \mathcal{H}_{\mathcal{F}}\right)  $
\begin{equation}
\mathcal{\Xi}_{\mathcal{H}_{\mathcal{F}}\mathcal{H}_{\mathcal{F}}^{d}}%
:\wedge\left(  \mathcal{H}^{1}\right)  \otimes\wedge\left(  \mathcal{H}%
^{1d}\right)  \rightarrow\mathcal{F}%
\end{equation}
is given by
\begin{equation}
\mathcal{\Xi}_{\mathcal{H}_{\mathcal{F}}\mathcal{H}_{\mathcal{F}}^{d}}\left(
\left\vert \Psi\right\rangle \otimes\left\langle \Phi\right\vert \right)
=\left\vert \Psi\right\rangle \left\langle \Phi\right\vert .
\end{equation}

\end{definition}

\begin{definition}
\begin{enumerate}
\item Multiplication "$\cdot$" in $\wedge\left(  \mathcal{H}^{1}\right)
\otimes\wedge\left(  \mathcal{H}^{1d}\right)  $ is defined by
\begin{equation}
\left\vert \Psi_{1}\right\rangle \otimes\left\langle \Phi_{1}\right\vert
\cdot\left\vert \Psi_{2}\right\rangle \otimes\left\langle \Phi_{2}\right\vert
=\left\vert \Psi_{1}\wedge\Psi_{2}\right\rangle \otimes\left\langle \Phi
_{1}\wedge\Phi_{2}\right\vert
\end{equation}


\item Multiplication "$\cdot$" in $\mathcal{F=B}\left(  \mathcal{H}%
_{\mathcal{F}}\right)  $ is defined by
\begin{equation}
\left\vert \Psi_{1}\right\rangle \left\langle \Phi_{1}\right\vert
\cdot\left\vert \Psi_{2}\right\rangle \left\langle \Phi_{2}\right\vert
=\left\vert \Psi_{1}\wedge\Psi_{2}\right\rangle \left\langle \Phi_{1}%
\wedge\Phi_{2}\right\vert
\end{equation}


\item The restriction \textquotedblright$\left.  \cdot\right\vert _{\left(
\wedge\left(  \mathcal{H}^{1}\right)  \otimes\wedge\left(  \mathcal{H}%
^{1d}\right)  \right)  _{e}}$\textquotedblright\ of \textquotedblright$\cdot
$\textquotedblright\ to the even spaces $\left(  \wedge\left(  \mathcal{H}%
^{1}\right)  \otimes\wedge\left(  \mathcal{H}^{1d}\right)  \right)  _{e}$ and
$\Delta\left(  \mathcal{H}^{1}\right)  $ are denoted by $\boxtimes.$
\end{enumerate}
\end{definition}

\begin{notation}
We use the same symbol in $\wedge\left(  \mathcal{H}^{1}\right)  \otimes
\wedge\left(  \mathcal{H}^{1d}\right)  $ and $\mathcal{B}\left(
\mathcal{H}_{\mathcal{F}}\right)  $ for multiplication.
\end{notation}

\begin{lemma}
$\mathcal{\Xi}_{\mathcal{H}_{\mathcal{F}}\mathcal{H}_{\mathcal{F}}^{d}}$ is an
algebraic map
\begin{align}
\mathcal{\Xi}_{\mathcal{H}_{\mathcal{F}}\mathcal{H}_{\mathcal{F}}^{d}}\left(
\left\vert \Psi_{1}\right\rangle \otimes\left\langle \Phi_{1}\right\vert
\cdot\left\vert \Psi_{2}\right\rangle \otimes\left\langle \Phi_{2}\right\vert
\right)   &  =\mathcal{\Xi}_{\mathcal{H}_{\mathcal{F}}\mathcal{H}%
_{\mathcal{F}}^{d}}\left(  \left\vert \Psi_{1}\wedge\Psi_{2}\right\rangle
\otimes\left\langle \Phi_{1}\wedge\Phi_{2}\right\vert \right) \\
&  =\mathcal{\Xi}_{\mathcal{H}_{\mathcal{F}}\mathcal{H}_{\mathcal{F}}^{d}%
}\left(  \left\vert \Psi_{1}\right\rangle \otimes\left\langle \Phi
_{1}\right\vert \right)  \cdot\mathcal{\Xi}_{\mathcal{H}_{\mathcal{F}%
}\mathcal{H}_{\mathcal{F}}^{d}}\left(  \left\vert \Psi_{2}\right\rangle
\otimes\left\langle \Phi_{2}\right\vert \right)  =\left\vert \Psi
_{1}\right\rangle \otimes\left\langle \Phi_{1}\right\vert \cdot\left(
\left\vert \Psi_{2}\right\rangle \otimes\left\langle \Phi_{2}\right\vert
\right)
\end{align}

\end{lemma}

\begin{corollary}
The restriction $\left.  \mathcal{\Xi}_{\mathcal{H}_{\mathcal{F}}%
\mathcal{H}_{\mathcal{F}}^{d}}\right\vert _{\left(  \wedge\left(
\mathcal{H}^{1}\right)  \otimes\wedge\left(  \mathcal{H}^{1d}\right)  \right)
_{e}}$ of $\mathcal{\Xi}_{\mathcal{H}_{\mathcal{F}}\mathcal{H}_{\mathcal{F}%
}^{d}}$ to the even subalgebra $\left(  \wedge\left(  \mathcal{H}^{1}\right)
\otimes\wedge\left(  \mathcal{H}^{1d}\right)  \right)  _{e}$ is an algebraic
map
\begin{align}
\mathcal{\Xi}_{\mathcal{H}_{\mathcal{F}}\mathcal{H}_{\mathcal{F}}^{d}}\left(
\left\vert \Psi_{1}\right\rangle \otimes\left\langle \Phi_{1}\right\vert
\boxtimes\left\vert \Psi_{2}\right\rangle \otimes\left\langle \Phi
_{2}\right\vert \right)   &  =\mathcal{\Xi}_{\mathcal{H}_{\mathcal{F}%
}\mathcal{H}_{\mathcal{F}}^{d}}\left(  \left\vert \Psi_{1}\wedge\Psi
_{2}\right\rangle \boxtimes\left\langle \Phi_{1}\wedge\Phi_{2}\right\vert
\right) \\
&  =\mathcal{\Xi}_{\mathcal{H}_{\mathcal{F}}\mathcal{H}_{\mathcal{F}}^{d}%
}\left(  \left\vert \Psi_{1}\right\rangle \otimes\left\langle \Phi
_{1}\right\vert \right)  \boxtimes\mathcal{\Xi}_{\mathcal{H}_{\mathcal{F}%
}\mathcal{H}_{\mathcal{F}}^{d}}\left(  \left\vert \Psi_{2}\right\rangle
\otimes\left\langle \Phi_{2}\right\vert \right)  =\left\vert \Psi
_{1}\right\rangle \left\langle \Phi_{1}\right\vert \boxtimes\left\vert
\Psi_{2}\right\rangle \left\langle \Phi_{2}\right\vert .
\end{align}

\end{corollary}

\begin{remark}
\bigskip In $\left(  \wedge\left(  \mathcal{H}^{1}\right)  \otimes
\wedge\left(  \mathcal{H}^{1d}\right)  \right)  _{e}$ and its image
$\mathcal{F}_{e}$
\begin{align}
w^{0\boxtimes}  &  =P_{\mathcal{H}^{0}}\\
w^{k\boxtimes}  &  =\frac{1}{k!}w\cdot\cdot\cdot w,\quad k\geq1.
\end{align}

\end{remark}

\begin{example}%
\begin{align}
\left(  \left\vert \Psi_{1}\right\rangle \left\langle \Psi_{2}\right\vert
+\left\vert \Psi_{2}\right\rangle \left\langle \Psi_{1}\right\vert \right)
^{2\boxtimes}  &  =\left(  \left\vert \Psi_{1}\right\rangle \left\langle
\Psi_{2}\right\vert +\left\vert \Psi_{2}\right\rangle \left\langle \Psi
_{1}\right\vert \right)  \boxtimes\left(  \left\vert \Psi_{1}\right\rangle
\left\langle \Psi_{2}\right\vert +\left\vert \Psi_{2}\right\rangle
\left\langle \Psi_{1}\right\vert \right) \\
&  =\left\vert \Psi_{1}\wedge\Psi_{2}\right\rangle \left\langle \Psi_{2}%
\wedge\Psi_{1}\right\vert +\left\vert \Psi_{2}\wedge\Psi_{1}\right\rangle
\left\langle \Psi_{1}\wedge\Psi_{2}\right\vert \\
&  =2\left\vert \Psi_{1}\wedge\Psi_{2}\right\rangle \left\langle \Psi
_{2}\wedge\Psi_{2}\right\vert =-2\left\vert \Psi_{1}\right\rangle \left\langle
\Psi_{1}\right\vert \boxtimes\left\vert \Psi_{2}\right\rangle \left\langle
\Psi_{2}\right\vert
\end{align}
and
\begin{align}
\left(  \left\vert \Psi_{1}\right\rangle \left\langle \Psi_{1}\right\vert
+\left\vert \Psi_{2}\right\rangle \left\langle \Psi_{2}\right\vert \right)
^{2\boxtimes}  &  =\left(  \left\vert \Psi_{1}\right\rangle \left\langle
\Psi_{1}\right\vert +\left\vert \Psi_{2}\right\rangle \left\langle \Psi
_{2}\right\vert \right)  \boxtimes\left(  \left\vert \Psi_{1}\right\rangle
\left\langle \Psi_{1}\right\vert +\left\vert \Psi_{2}\right\rangle
\left\langle \Psi_{2}\right\vert \right) \\
&  =\left\vert \Psi_{1}\wedge\Psi_{2}\right\rangle \left\langle \Psi_{1}%
\wedge\Psi_{2}\right\vert +\left\vert \Psi_{2}\wedge\Psi_{1}\right\rangle
\left\langle \Psi_{2}\wedge\Psi_{1}\right\vert \\
&  =2\left\vert \Psi_{1}\wedge\Psi_{2}\right\rangle \left\langle \Psi
_{1}\wedge\Psi_{2}\right\vert =2\left\vert \Psi_{1}\right\rangle \left\langle
\Psi_{1}\right\vert \boxtimes\left\vert \Psi_{2}\right\rangle \left\langle
\Psi_{2}\right\vert .
\end{align}

\end{example}

\begin{remark}%
\begin{equation}
2\left\vert \Psi_{1}\wedge\Psi_{2}\right\rangle \left\langle \Psi_{1}%
\wedge\Psi_{2}\right\vert =\left\vert \Psi_{1}\Psi_{2}\right\rangle
\left\langle \Psi_{1}\Psi_{2}\right\vert ,
\end{equation}
where $\left\vert \Psi_{1}\Psi_{2}\right\rangle $ are normalized and
\begin{align}
\operatorname{Tr}\left\{  \left\vert \Psi_{1}\Psi_{2}\right\rangle
\left\langle \Psi_{1}\Psi_{2}\right\vert \right\}   &  =1\\
\left\vert \Psi_{1}\Psi_{2}\right\rangle \left\langle \Psi_{1}\Psi
_{2}\right\vert ^{2}  &  =\left\vert \Psi_{1}\Psi_{2}\right\rangle
\left\langle \Psi_{1}\Psi_{2}\right\vert .
\end{align}

\end{remark}

The mixed exterior algebra $%
%TCIMACRO{\tbigwedge }%
%BeginExpansion
{\textstyle\bigwedge}
%EndExpansion
\left(  \mathcal{H}^{1d},\mathcal{H}^{1}\right)  $ is \emph{bigraded} as%
\begin{equation}%
%TCIMACRO{\tbigwedge }%
%BeginExpansion
{\textstyle\bigwedge}
%EndExpansion
\left(  \mathcal{H}^{1},\mathcal{H}^{1d}\right)  =%
%TCIMACRO{\tbigoplus \limits_{0\leq p,q\leq r}}%
%BeginExpansion
{\textstyle\bigoplus\limits_{0\leq p,q\leq r}}
%EndExpansion%
%TCIMACRO{\tbigwedge \nolimits_{p}^{q}}%
%BeginExpansion
{\textstyle\bigwedge\nolimits_{p}^{q}}
%EndExpansion
\left(  \mathcal{H}^{1},\mathcal{H}^{1d}\right)  ,
\end{equation}
where
\begin{equation}%
%TCIMACRO{\tbigwedge \nolimits_{p}^{q}}%
%BeginExpansion
{\textstyle\bigwedge\nolimits_{p}^{q}}
%EndExpansion
\left(  \mathcal{H}^{1},\mathcal{H}^{1d}\right)  =%
%TCIMACRO{\tbigwedge \nolimits^{p}}%
%BeginExpansion
{\textstyle\bigwedge\nolimits^{p}}
%EndExpansion
\left(  \mathcal{H}^{1}\right)  \otimes%
%TCIMACRO{\tbigwedge \nolimits^{q}}%
%BeginExpansion
{\textstyle\bigwedge\nolimits^{q}}
%EndExpansion
\left(  \mathcal{H}^{1d}\right)  .
\end{equation}
Clearly
\begin{equation}
w_{1}\cdot w_{2}=\left(  -1\right)  ^{p_{1}p_{2}+q_{1}q_{2}}w_{2}\cdot
w_{1};\quad w_{1}\in%
%TCIMACRO{\tbigwedge \nolimits_{q_{1}}^{p_{1}}}%
%BeginExpansion
{\textstyle\bigwedge\nolimits_{q_{1}}^{p_{1}}}
%EndExpansion
\left(  \mathcal{H}^{1},\mathcal{H}^{1d}\right)  ;w_{2}\in%
%TCIMACRO{\tbigwedge \nolimits_{q_{1}}^{p_{2}}}%
%BeginExpansion
{\textstyle\bigwedge\nolimits_{q_{1}}^{p_{2}}}
%EndExpansion
\left(  \mathcal{H}^{1},\mathcal{H}^{1d}\right)  .
\end{equation}
This gives
\begin{equation}
w_{1}\cdot w_{2}=w_{2}\cdot w_{1}\text{ if }p_{1}+q_{1}\text{ and }p_{2}%
+q_{2}\text{ are even}%
\end{equation}
and
\begin{equation}
w_{1}\cdot w_{2}=-w_{2}\cdot w_{1}\text{ if }p_{1}+q_{1}\text{ and }%
p_{2}+q_{2}\text{ are not both even.}%
\end{equation}
In the case that both are even one has the binomial formula
\begin{equation}
\left(  w_{1}+w_{2}\right)  ^{l}=\sum_{j+k=l}w_{1}^{j}\cdot w_{2}^{k}.
\end{equation}
A \emph{scalar product (not necessarily inner) in }$%
%TCIMACRO{\tbigwedge \nolimits^{p}}%
%BeginExpansion
{\textstyle\bigwedge\nolimits^{p}}
%EndExpansion
\left(  \mathcal{H}^{1}\right)  \times%
%TCIMACRO{\tbigwedge \nolimits^{q}}%
%BeginExpansion
{\textstyle\bigwedge\nolimits^{q}}
%EndExpansion
\left(  \mathcal{H}^{1d}\right)  .$\emph{ is defined by }%
\begin{equation}
\left\langle x_{1}^{d}\wedge\cdots\wedge x_{p}^{d},x_{1}\wedge\cdots\wedge
x_{q}\right\rangle =\left\{
\begin{array}
[c]{c}%
\det\left\{  \left\langle x_{j}^{d},x_{k}\right\rangle \right\}  \text{ if
}p=q\\
0\text{ if }p\neq q
\end{array}
\right.
\end{equation}
induces an \emph{inner product} in $%
%TCIMACRO{\tbigwedge }%
%BeginExpansion
{\textstyle\bigwedge}
%EndExpansion
\left(  \mathcal{H}^{1},\mathcal{H}^{1d}\right)  $ via
\begin{equation}
\left\langle \left.  u\otimes v^{d}\right\vert w\otimes x^{d}\right\rangle
=\left\langle v^{d},w\right\rangle \left\langle x^{d},u\right\rangle \text{ if
}\mathcal{H}^{1}\text{ is a Hilbert space then}\overset{\text{ }%
\mathcal{H}^{1}\overset{anti}{\cong}\mathcal{H}^{1d}}{=}\left\langle \left.
v\right\vert w\right\rangle \left\langle \left.  x\right\vert u\right\rangle
\end{equation}
that corresponds to%
\begin{align}
&  \left\langle \left.  \Psi_{1}\otimes\Phi_{1}^{\dagger}\right\vert \Psi
_{2}\otimes\Phi_{2}^{\dagger}\right\rangle =\left\langle \Phi_{1}^{\dagger
},\Psi_{2}\right\rangle \left\langle \Phi_{2}^{\dagger},\Psi_{1}\right\rangle
\overset{\text{ }\mathcal{H}\cong\mathcal{H}^{d}}{=}\left\langle \left.
\Phi_{2}\right\vert \Psi_{1}\right\rangle \left\langle \left.  \Phi
_{1}\right\vert \Psi_{2}\right\rangle \\
&  =\operatorname{Tr}\left\{  \left(  \left\vert \Phi_{1}\right\rangle
\left\langle \Psi_{1}\right\vert \right)  ^{\dagger}\left\vert \Psi
_{2}\right\rangle \left\langle \Phi_{2}\right\vert \right\}
=\operatorname{Tr}\left\{  \left\vert \Psi_{1}\right\rangle \left\langle
\Phi_{1}\right\vert \left\vert \Psi_{2}\right\rangle \left\langle \Phi
_{2}\right\vert \right\}  =\left\langle \left.  \overset{}{\left\vert \Phi
_{1}\right\rangle \left\langle \Psi_{1}\right\vert }\right\vert \left\vert
\Psi_{2}\right\rangle \left\langle \Phi_{2}\right\vert \right\rangle ,
\end{align}
which is an inner product.

\begin{remark}
This product is inner even if there is no inner product in $\mathcal{H}^{1}.$
\end{remark}

\subsubsection{Creation and Annihilation Operators in the Exterior Algebra $%
%TCIMACRO{\tbigwedge }%
%BeginExpansion
{\textstyle\protect\bigwedge}
%EndExpansion
\left(  \mathcal{H}^{1}\right)  $}

We will asume that $\mathcal{H}^{1}$ is a Hilbert space unless specified
otherwise. The multiplication operator $\mu\left(  z\right)  $ on the exterior
algebra $%
%TCIMACRO{\tbigwedge }%
%BeginExpansion
{\textstyle\bigwedge}
%EndExpansion
\left(  \mathcal{H}\right)  $ is defined by
\begin{align}
\mu\left(  z\right)   &  :%
%TCIMACRO{\tbigwedge }%
%BeginExpansion
{\textstyle\bigwedge}
%EndExpansion
\left(  \mathcal{H}^{1}\right)  \rightarrow%
%TCIMACRO{\tbigwedge }%
%BeginExpansion
{\textstyle\bigwedge}
%EndExpansion
\left(  \mathcal{H}^{1}\right) \\
\mu\left(  z\right)  w  &  =z\wedge w;\quad z,w\in%
%TCIMACRO{\tbigwedge }%
%BeginExpansion
{\textstyle\bigwedge}
%EndExpansion
\left(  \mathcal{H}^{1}\right)
\end{align}
and the dual operator $i\left(  z\right)  $%
\begin{equation}
i\left(  z\right)  ^{d}:%
%TCIMACRO{\tbigwedge }%
%BeginExpansion
{\textstyle\bigwedge}
%EndExpansion
\left(  \mathcal{H}^{1d}\right)  \rightarrow%
%TCIMACRO{\tbigwedge }%
%BeginExpansion
{\textstyle\bigwedge}
%EndExpansion
\left(  \mathcal{H}^{1d}\right)  ;\quad z\in%
%TCIMACRO{\tbigwedge }%
%BeginExpansion
{\textstyle\bigwedge}
%EndExpansion
\left(  \mathcal{H}^{1}\right)
\end{equation}
by
\begin{equation}
\left\langle i\left(  z\right)  ^{d}w_{1}^{d},w_{2}\right\rangle =\left\langle
w_{1}^{d},\mu\left(  z\right)  w_{2}\right\rangle ;\quad z,w_{2}\in%
%TCIMACRO{\tbigwedge }%
%BeginExpansion
{\textstyle\bigwedge}
%EndExpansion
\left(  \mathcal{H}\right)  ,w_{1}^{d},\in%
%TCIMACRO{\tbigwedge }%
%BeginExpansion
{\textstyle\bigwedge}
%EndExpansion
\left(  \mathcal{H}^{d}\right)  .
\end{equation}
One can show that
\begin{equation}
\mu\left(  z_{1}\wedge z_{2}\right)  =\mu\left(  z_{1}\right)  \circ\mu\left(
z_{2}\right)  ;\quad z_{1},z_{2}\in%
%TCIMACRO{\tbigwedge }%
%BeginExpansion
{\textstyle\bigwedge}
%EndExpansion
\left(  \mathcal{H}^{1}\right)  .
\end{equation}
and
\begin{equation}
i\left(  z_{1}\wedge z_{2}\right)  ^{d}=i\left(  z_{2}\right)  ^{d}\circ
i\left(  z_{1}\right)  ^{d};\quad z_{1},z_{2}\in%
%TCIMACRO{\tbigwedge }%
%BeginExpansion
{\textstyle\bigwedge}
%EndExpansion
\left(  \mathcal{H}^{1}\right)  ,
\end{equation}
thus one can build any $\mu\left(  z\right)  $ and $i\left(  z\right)  ^{d}$
from $\left\{  \left.  \mu\left(  z\right)  ,i\left(  z\right)  ^{d}%
\right\vert z\in\mathcal{H}^{1}\right\}  .$

In particular when $z\in\mathcal{H}^{1}$ the effect of $i\left(  z\right)  $
is given by
\begin{equation}
i\left(  z\right)  ^{d}\left(  y_{k_{1}}^{d}\wedge\cdots\wedge y_{k_{p}}%
^{d}\right)  =\sum_{1\leq\mu\leq p}\left(  -1\right)  ^{\mu-1}\left\langle
y_{k_{\mu}}^{d},z\right\rangle y_{k_{1}}^{d}\wedge\cdots\wedge
\widehat{y_{k_{\mu}}^{d}}\wedge\cdots\wedge y_{k_{p}}^{d},
\end{equation}
which can be expressed in the Hilbert space context, after using the
antilinear isometric isomorphism $\dagger$, as
\begin{equation}
i\left(  z\right)  ^{d}\left(  \left\langle y_{k_{1}}\right\vert \wedge
\cdots\wedge\left\langle y_{k_{p}}\right\vert \right)  =\sum_{1\leq\mu\leq
p}\left(  -1\right)  ^{\mu-1}\left\langle y_{k_{\mu}}\left\vert z\right.
\right\rangle \left\langle y_{k_{1}}\right\vert \wedge\cdots\wedge\left\langle
\widehat{y_{k_{\mu}}}\right\vert \wedge\cdots\wedge\left\langle y_{k_{p}%
}\right\vert .
\end{equation}


E.g.%
\begin{equation}
i\left(  z\right)  ^{d}\left(  \left\langle y_{1}\right\vert \wedge
\left\langle y_{2}\right\vert \right)  =\left(  \left\langle y_{1}\right\vert
\wedge\left\langle y_{2}\right\vert \right)  \mu\left(  z\right)
=\left\langle y_{1}\left\vert z\right.  \right\rangle \left\langle
y_{2}\right\vert -\left\langle y_{2}\left\vert z\right.  \right\rangle
\left\langle y_{1}\right\vert ,
\end{equation}
which under the "$\dagger$" isomorphism gives
\begin{equation}
\left(  \left(  \left\langle y_{1}\right\vert \wedge\left\langle
y_{2}\right\vert \right)  \mu\left(  z\right)  \right)  ^{\dagger}=\mu\left(
z\right)  ^{\dagger}\left\vert y_{1}\right\rangle \wedge\left\vert
y_{2}\right\rangle =i\left(  z\right)  \left\vert y_{1}\right\rangle
\wedge\left\vert y_{2}\right\rangle =\left\langle z\left\vert y_{1}\right.
\right\rangle \left\langle y_{2}\right\vert -\left\langle y_{1}\right\vert
\left\langle z\left\vert y_{2}\right.  \right\rangle .
\end{equation}


\subsubsection{Creation and Annihilation Operators in the Mixed Exterior
Algebra $%
%TCIMACRO{\tbigwedge }%
%BeginExpansion
{\textstyle\protect\bigwedge}
%EndExpansion
\left(  \mathcal{H}^{1},\mathcal{H}^{1d}\right)  $}

The multiplication operator $\mu\left(  z\right)  $ on the mixed algebra is
defined by
\begin{equation}
\mu\left(  z\right)  w=z\cdot w;\quad z,w\in%
%TCIMACRO{\tbigwedge }%
%BeginExpansion
{\textstyle\bigwedge}
%EndExpansion
\left(  \mathcal{H}^{1},\mathcal{H}^{1d}\right)  .
\end{equation}


\begin{example}
Let
\begin{equation}
z=\left\vert x_{j_{1}}\right\rangle \left\langle x_{k_{1}}\right\vert
\end{equation}
then%
\begin{gather}
\mu\left(  \left\vert x_{j_{1}}\right\rangle \left\langle x_{k_{1}}\right\vert
\right)  \left\vert x_{l_{1}}\wedge\cdots\wedge x_{l_{p}}\right\rangle
\left\langle x_{m_{1}}\wedge\cdots\wedge x_{m_{p}}\right\vert =\mu\left(
\left\vert x_{j_{1}}\right\rangle \right)  \left\vert x_{l_{1}}\wedge
\cdots\wedge x_{l_{p}}\right\rangle \left\langle x_{m_{1}}\wedge\cdots\wedge
x_{m_{p}}\right\vert \left(  \mu\left(  \left\langle x_{k_{1}}\right\vert
\right)  \right)  ^{\dagger}\\
=\left\vert \mu\left(  \left\vert x_{j_{1}}\right\rangle \right)  \left(
x_{l_{1}}\wedge\cdots\wedge x_{l_{p}}\right)  \right\rangle \left\langle
\mu\left(  \left\vert x_{k_{1}}\right\rangle \right)  \left(  x_{m_{1}}%
\wedge\cdots\wedge x_{m_{p}}\right)  \right\vert =\left\vert x_{j_{1}}\wedge
x_{l_{1}}\wedge\cdots\wedge x_{l_{p}}\right\rangle \left\langle x_{k_{1}%
}\wedge x_{m_{1}}\wedge\cdots\wedge x_{m_{p}}\right\vert .
\end{gather}

\end{example}

\bigskip The dual operator $i\left(  z\right)  $ is given by
\begin{equation}
\left\langle \left.  i\left(  z\right)  w_{1}\right\vert w_{2}\right\rangle
=\left\langle \left.  w_{1}\right\vert \mu\left(  z\right)  w_{2}\right\rangle
;\quad w_{1},w_{2}\in%
%TCIMACRO{\tbigwedge }%
%BeginExpansion
{\textstyle\bigwedge}
%EndExpansion
\left(  \mathcal{H}^{1},\mathcal{H}^{1d}\right)  .
\end{equation}


\begin{example}
Let
\begin{equation}
z=\left\vert x_{j_{1}}\right\rangle \left\langle x_{k_{1}}\right\vert
\end{equation}
then
\begin{align}
&  i\left(  \left\vert x_{j_{1}}\right\rangle \left\langle x_{k_{1}%
}\right\vert \right)  \left(  \left\vert x_{l_{1}}\wedge\cdots\wedge x_{l_{p}%
}\right\rangle \left\langle x_{m_{1}}\wedge\cdots\wedge x_{m_{p}}\right\vert
\right)  =i\left(  \left\langle x_{k_{1}}\right\vert \right)  \left(
\left\vert x_{l_{1}}\wedge\cdots\wedge x_{l_{p}}\right\rangle \left\langle
x_{m_{1}}\wedge\cdots\wedge x_{m_{p}}\right\vert \right)  i\left(  \left\vert
x_{j_{1}}\right\rangle \right)  ^{\dagger}\\
&  =\left\vert i\left(  \left\langle x_{k_{1}}\right\vert \right)  \left(
x_{l_{1}}\wedge\cdots\wedge x_{l_{p}}\right)  \right\rangle \left\langle
i\left(  \left\vert x_{j_{1}}\right\rangle \right)  \left(  x_{m_{1}}%
\wedge\cdots\wedge x_{m_{p}}\right)  \right\vert \\
&  =\sum_{1\leq v,\mu\leq p}\left(  -1\right)  ^{v-1}\left\langle \left.
x_{k_{1}}\right\vert x_{l_{\nu}}\right\rangle \left\vert x_{_{l_{1}}}%
\wedge\cdots\wedge\widehat{x_{l_{\nu}}}\wedge\cdots\wedge x_{l_{p}%
}\right\rangle \left\langle x_{m_{1}}\wedge\cdots\wedge\widehat{x_{m_{\mu}}%
}\wedge\cdots\wedge x_{m_{p}}\right\vert \left(  -1\right)  ^{\mu
-1}\left\langle \left.  x_{m_{\mu}}\right\vert x_{j_{1}}\right\rangle \\
&  =\sum_{1\leq v,\mu\leq p}\left(  -1\right)  ^{\mu+v}\left\langle \left.
x_{m_{\mu}}\right\vert x_{j_{1}}\right\rangle \left\langle \left.  x_{k_{1}%
}\right\vert x_{l_{\nu}}\right\rangle \left\vert x_{_{l_{1}}}\wedge
\cdots\wedge\widehat{x_{l_{\nu}}}\wedge\cdots\wedge x_{l_{p}}\right\rangle
\left\langle x_{m_{1}}\wedge\cdots\wedge\widehat{x_{m_{\mu}}}\wedge
\cdots\wedge x_{m_{p}}\right\vert .
\end{align}
$\lceil$The $i\left(  \left\vert x_{2}\right\rangle \left\langle
x_{1}\right\vert \right)  $ operator on $%
%TCIMACRO{\tbigwedge _{2}^{2}}%
%BeginExpansion
{\textstyle\bigwedge_{2}^{2}}
%EndExpansion
\left(  \mathcal{H}^{1},\mathcal{H}^{1d}\right)  $ gives
\begin{equation}
i\left(  \overset{}{\left\vert x_{2}\right\rangle \left\langle x_{1}%
\right\vert }\right)  \left(  \overset{}{\left\vert x_{1}\wedge x_{3}%
\right\rangle \left\langle x_{2}\wedge x_{4}\right\vert }\right)  =\left(
\left\vert i\left(  \left\langle x_{1}\right\vert \right)  x_{1}\wedge
x_{3}\right\rangle \left\langle i\left(  \left\langle x_{2}\right\vert
\right)  x_{2}\wedge x_{4}\right\vert \right)  =\left\vert x_{3}\right\rangle
\left\langle x_{4}\right\vert
\end{equation}
$\rfloor$
\end{example}

Hence
\begin{equation}
\mu\left(  \left\vert x_{j_{1}}\wedge\cdots\wedge x_{j_{p}}\right\rangle
\left\langle x_{k_{1}}\wedge\cdots\wedge x_{k_{q}}\right\vert \right)  :%
%TCIMACRO{\tbigwedge \nolimits_{n}^{m}}%
%BeginExpansion
{\textstyle\bigwedge\nolimits_{n}^{m}}
%EndExpansion
\left(  \mathcal{H}^{1},\mathcal{H}^{1d}\right)  \rightarrow%
%TCIMACRO{\tbigwedge \nolimits_{n+p}^{m+q}}%
%BeginExpansion
{\textstyle\bigwedge\nolimits_{n+p}^{m+q}}
%EndExpansion
\left(  \mathcal{H}^{1},\mathcal{H}^{1d}\right)
\end{equation}
and%
\begin{equation}
i\left(  \left\vert x_{j_{1}}\wedge\cdots\wedge x_{j_{p}}\right\rangle
\left\langle x_{k_{1}}\wedge\cdots\wedge x_{k_{q}}\right\vert \right)  :%
%TCIMACRO{\tbigwedge \nolimits_{n}^{m}}%
%BeginExpansion
{\textstyle\bigwedge\nolimits_{n}^{m}}
%EndExpansion
\left(  \mathcal{H}^{1},\mathcal{H}^{1d}\right)  \rightarrow%
%TCIMACRO{\tbigwedge \nolimits_{n-p}^{m-q}}%
%BeginExpansion
{\textstyle\bigwedge\nolimits_{n-p}^{m-q}}
%EndExpansion
\left(  \mathcal{H}^{1},\mathcal{H}^{1d}\right)  .
\end{equation}


It follows from the definition that
\begin{equation}
\mu\left(  w_{1}\otimes w_{2}^{d}\right)  =\mu\left(  w_{1}\right)  \otimes
\mu\left(  w_{2}^{d}\right)  ;\quad w_{1},w_{2}\in%
%TCIMACRO{\tbigwedge }%
%BeginExpansion
{\textstyle\bigwedge}
%EndExpansion
\left(  \mathcal{H},\mathcal{H}^{d}\right)  .
\end{equation}
In particular
\begin{align}
\mu\left(  I_{\mathcal{H}}\otimes w^{d}\right)   &  =\mu\left(  I_{\mathcal{H}%
}\right)  \otimes\mu\left(  w^{d}\right)  =I_{\mathcal{H}}\otimes\mu\left(
w^{d}\right) \\
\mu\left(  w\otimes I_{\mathcal{H}^{d}}\right)   &  =\mu\left(  w\right)
\otimes\mu\left(  I_{\mathcal{H}^{d}}\right)  =w\otimes I_{\mathcal{H}^{d}}.
\end{align}
Dualizing the relation
\begin{equation}
\mu\left(  z_{1}\cdot z_{2}\right)  =\mu\left(  z_{1}\right)  \circ\mu\left(
z_{2}\right)  ;\quad z_{1},z_{2}\in%
%TCIMACRO{\tbigwedge }%
%BeginExpansion
{\textstyle\bigwedge}
%EndExpansion
\left(  \mathcal{H}^{1},\mathcal{H}^{1d}\right)
\end{equation}
one obtains
\begin{equation}
i\left(  z_{1}\cdot z_{2}\right)  =i\left(  z_{2}\right)  \circ i\left(
z_{1}\right)  ;\quad z_{1},z_{2}\in%
%TCIMACRO{\tbigwedge }%
%BeginExpansion
{\textstyle\bigwedge}
%EndExpansion
\left(  \mathcal{H}^{1},\mathcal{H}^{1d}\right)  .
\end{equation}
Note that if $z\in%
%TCIMACRO{\tbigwedge \nolimits_{q}^{p}}%
%BeginExpansion
{\textstyle\bigwedge\nolimits_{q}^{p}}
%EndExpansion
\left(  \mathcal{H}^{1},\mathcal{H}^{1d}\right)  $ and $w\in%
%TCIMACRO{\tbigwedge \nolimits_{p}^{q}}%
%BeginExpansion
{\textstyle\bigwedge\nolimits_{p}^{q}}
%EndExpansion
\left(  \mathcal{H}^{1},\mathcal{H}^{1d}\right)  $ then $i\left(  z\right)
w\in%
%TCIMACRO{\tbigwedge \nolimits_{0}^{0}}%
%BeginExpansion
{\textstyle\bigwedge\nolimits_{0}^{0}}
%EndExpansion
\left(  \mathcal{H}^{1},\mathcal{H}^{1d}\right)  $ given by
\begin{equation}
i\left(  z\right)  w=\left\langle \left.  z\right\vert w\right\rangle .
\end{equation}


\subsubsection{The Diagonal Subalgebra}

The Diagonal subsalgebra is the commutative subalgebra defined by
\begin{equation}
\Delta\left(  \mathcal{H}^{1},\mathcal{H}^{1d}\right)  =%
%TCIMACRO{\tbigoplus \limits_{0\leq p\leq n}}%
%BeginExpansion
{\textstyle\bigoplus\limits_{0\leq p\leq n}}
%EndExpansion
\Delta_{p}\left(  \mathcal{H}^{1},\mathcal{H}^{1d}\right)  ;\quad\Delta
_{p}\left(  \mathcal{H}^{1}\right)  =%
%TCIMACRO{\tbigwedge \nolimits_{p}^{p}}%
%BeginExpansion
{\textstyle\bigwedge\nolimits_{p}^{p}}
%EndExpansion
\left(  \mathcal{H}^{1},\mathcal{H}^{1d}\right)  .
\end{equation}


\begin{enumerate}
\item The diagonal subalgebra is stable under $i\left(  z\right)  $ if
$z\in\Delta\left(  \mathcal{H}^{1},\mathcal{H}^{1d}\right)  .$

\item The restriction of the inner product to $\Delta\left(  \mathcal{H}%
^{1},\mathcal{H}^{1d}\right)  $ is nondegenerate.
\end{enumerate}

\subsubsection{The Isomorphism $\Delta\left(  \mathcal{H}^{1},\mathcal{H}%
^{1d}\right)  \protect\cong\mathcal{B}\left(  \mathcal{H}_{\mathcal{F}%
}\right)  _{e}$}

The isomorphism
\begin{equation}
\Xi_{\mathcal{H}_{\mathcal{F}}\mathcal{H}_{\mathcal{F}}^{d}}:\Delta\left(
\mathcal{H}^{1},\mathcal{H}^{1d}\right)  \rightarrow\mathcal{B}\left(
\mathcal{H}_{\mathcal{F}}\right)  _{e}%
\end{equation}
is defined by
\begin{equation}
\Xi_{\mathcal{H}_{\mathcal{F}}\mathcal{H}_{\mathcal{F}}^{d}}\left(  b\otimes
a^{d}\right)  v=\left\langle a^{d},v\right\rangle b=\left\langle \left.
a\right\vert v\right\rangle b.
\end{equation}
$\lceil$%
\begin{equation}
\Xi_{\mathcal{H}_{\mathcal{F}}\mathcal{H}_{\mathcal{F}}^{d}}\left(  \left\vert
\Psi_{1}\right\rangle \otimes\left\langle \Phi_{1}\right\vert \right)
\left\vert \Psi_{2}\right\rangle =\left\vert \Psi_{1}\right\rangle
\left\langle \Phi_{1}\right\vert \left\vert \Psi_{2}\right\rangle
=\left\langle \left.  \Phi_{1}\right\vert \Psi_{2}\right\rangle \left\vert
\Psi_{1}\right\rangle .
\end{equation}
$\rfloor$

Since
\begin{equation}
\left\langle \left.  a\right\vert v\right\rangle =0\text{ for }a^{d}\in%
%TCIMACRO{\tbigwedge \nolimits^{p}}%
%BeginExpansion
{\textstyle\bigwedge\nolimits^{p}}
%EndExpansion
\left(  \mathcal{H}^{1d}\right)  ,v\in%
%TCIMACRO{\tbigwedge \nolimits^{q}}%
%BeginExpansion
{\textstyle\bigwedge\nolimits^{q}}
%EndExpansion
\left(  \mathcal{H}^{1}\right)  ,p\neq q,
\end{equation}
$\Xi_{\mathcal{H}_{\mathcal{F}}\mathcal{H}_{\mathcal{F}}^{d}}$ restricts to
linear maps
\begin{equation}
\Xi_{\mathcal{H}_{\mathcal{F}}\mathcal{H}_{\mathcal{F}}^{d}}:%
%TCIMACRO{\tbigwedge \nolimits_{q}^{p}}%
%BeginExpansion
{\textstyle\bigwedge\nolimits_{q}^{p}}
%EndExpansion
\left(  \mathcal{H}^{1},\mathcal{H}^{1d}\right)  \rightarrow\mathcal{B}\left(
\mathcal{H}^{p},\mathcal{H}^{q}\right)  .
\end{equation}


The image of the diagonal subalgebra can be decomposed as
\begin{equation}
\Xi_{\mathcal{H}_{\mathcal{F}}\mathcal{H}_{\mathcal{F}}^{d}}\left(
\Delta\left(  \mathcal{H}^{1},\mathcal{H}^{1d}\right)  \right)  =\sum_{0\leq
n\leq r}\Xi_{\mathcal{H}_{\mathcal{F}}\mathcal{H}_{\mathcal{F}}^{d}}\left(
\Delta_{n}\left(  \mathcal{H}^{1}\right)  \right)  =%
%TCIMACRO{\tbigoplus \limits_{0\leq n\leq r}}%
%BeginExpansion
{\textstyle\bigoplus\limits_{0\leq n\leq r}}
%EndExpansion
\mathcal{B}\left(  \mathcal{H}^{n}\right)  .
\end{equation}


\subsubsection{Box Product of Linear Transformations $\in\mathcal{B}\left(
\mathcal{H}^{1}\right)  $}

\begin{definition}
The Box Product
\begin{equation}
\varphi_{1}\boxtimes\cdots\boxtimes\varphi_{p}:\mathcal{H}^{p}\rightarrow
\mathcal{H}^{p}%
\end{equation}
is defined by
\begin{equation}
\left(  \varphi_{1}\boxtimes\cdots\boxtimes\varphi_{p}\right)  \left(
x_{1}\wedge\cdots\wedge x_{p}\right)  =\frac{1}{p!}\sum_{\sigma\in S_{p}%
}\varphi_{1}\left(  x_{\sigma\left(  1\right)  }\right)  \wedge\cdots
\wedge\varphi_{p}\left(  x_{\sigma\left(  p\right)  }\right)  ,
\end{equation}
so that $\left(  \varphi_{1}\boxtimes\cdots\boxtimes\varphi_{p}\right)
\in\mathcal{B}\left(  \mathcal{H}^{p}\right)  .$
\end{definition}

One can observe that
\begin{equation}
\frac{1}{p!}\sum_{\sigma\in S_{p}}\varphi_{1}\left(  x_{\sigma\left(
1\right)  }\right)  \wedge\cdots\wedge\varphi_{p}\left(  x_{\sigma\left(
p\right)  }\right)  =\frac{1}{p!}\sum_{\sigma\in S_{p}}\varphi_{\sigma\left(
1\right)  }\left(  x_{1}\right)  \wedge\cdots\wedge\varphi_{\sigma\left(
p\right)  }\left(  x_{p}\right)  ,
\end{equation}
which shows that the box product is symmetric.

\begin{lemma}
The box product is a special case of the \textquotedblright$\cdot
$\textquotedblright\ product i.e.
\begin{equation}
\Xi_{\mathcal{H}_{\mathcal{F}}\mathcal{H}_{\mathcal{F}}^{d}}\left(  z_{1}%
\cdot\cdots\cdot z_{p}\right)  =\Xi_{\mathcal{H}_{\mathcal{F}}\mathcal{H}%
_{\mathcal{F}}^{d}}\left(  z_{1}\right)  \cdot\cdots\cdot\Xi_{\mathcal{H}%
_{\mathcal{F}}\mathcal{H}_{\mathcal{F}}^{d}}\left(  z_{p}\right)
=\Xi_{\mathcal{H}_{\mathcal{F}}\mathcal{H}_{\mathcal{F}}^{d}}\left(
z_{1}\right)  \boxtimes\cdots\boxtimes\Xi_{\mathcal{H}_{\mathcal{F}%
}\mathcal{H}_{\mathcal{F}}^{d}}\left(  z_{p}\right)  ;\quad z_{j}\in%
%TCIMACRO{\tbigwedge \nolimits_{1}^{1}}%
%BeginExpansion
{\textstyle\bigwedge\nolimits_{1}^{1}}
%EndExpansion
\left(  \mathcal{H}^{1},\mathcal{H}^{1d}\right)
\end{equation}

\end{lemma}

It follows from the definition of the box product that
\begin{equation}
\varphi\boxtimes\cdots\boxtimes\varphi=%
%TCIMACRO{\tbigwedge \nolimits^{p}}%
%BeginExpansion
{\textstyle\bigwedge\nolimits^{p}}
%EndExpansion
\varphi=%
%TCIMACRO{\tbigwedge \nolimits^{p}}%
%BeginExpansion
{\textstyle\bigwedge\nolimits^{p}}
%EndExpansion
\Xi_{\mathcal{H}_{\mathcal{F}}\mathcal{H}_{\mathcal{F}}^{d}}\left(  z\right)
=\Xi_{\mathcal{H}_{\mathcal{F}}\mathcal{H}_{\mathcal{F}}^{d}}\left(
z^{p}\right)  ;\quad z\in%
%TCIMACRO{\tbigwedge \nolimits_{1}^{1}}%
%BeginExpansion
{\textstyle\bigwedge\nolimits_{1}^{1}}
%EndExpansion
\left(  \mathcal{H}^{1},\mathcal{H}^{1d}\right)  .
\end{equation}


\begin{remark}
One should note that
\begin{equation}%
%TCIMACRO{\tbigwedge \nolimits^{p}}%
%BeginExpansion
{\textstyle\bigwedge\nolimits^{p}}
%EndExpansion
\varphi\neq\,\overset{\longleftarrow p\longrightarrow}{\varphi\wedge
\cdots\wedge\varphi}.
\end{equation}

\end{remark}

\subsubsection{Composition Product}

The composition product defined by
\begin{equation}
\left(  u_{1}\otimes u_{2}^{d}\right)  \circ\left(  v_{1}\otimes v_{2}%
^{d}\right)  =\left\langle u_{2}^{d},v_{1}\right\rangle u_{1}\otimes v_{2}^{d}%
\end{equation}
$\lceil$%
\begin{equation}
\left\vert \Psi_{1}\right\rangle \left\langle \Phi_{1}\right\vert
\cdot\left\vert \Psi_{2}\right\rangle \left\langle \Phi_{2}\right\vert
=\left\vert \Psi_{1}\wedge\Psi_{2}\right\rangle \left\langle \Phi_{1}%
\wedge\Phi_{2}\right\vert
\end{equation}%
\begin{equation}
\left\vert \Psi_{1}\right\rangle \left\langle \Phi_{1}\right\vert
\circ\left\vert \Psi_{2}\right\rangle \left\langle \Phi_{2}\right\vert
=\left\langle \Phi_{1}^{d},\Psi_{2}\right\rangle \left\vert \Psi
_{1}\right\rangle \left\langle \Phi_{2}\right\vert =\left\vert \Psi
_{1}\right\rangle \left\langle \Phi_{1}\right\vert \left\vert \Psi
_{2}\right\rangle \left\langle \Phi_{2}\right\vert =\left\langle \left.
\Phi_{1}\right\vert \Psi_{2}\right\rangle \left\vert \Psi_{1}\right\rangle
\left\langle \Phi_{2}\right\vert
\end{equation}
$\rfloor$

is the normal operator product in $\mathcal{B}\left(  \mathcal{H}%
_{\mathcal{F}}\right)  .$

\begin{lemma}
Let $\eta\in\mathcal{H}^{1}$
\begin{align}
i\left(  \eta\right)  \left(  y_{1}^{d}\wedge\cdots\wedge y_{p}^{d}\right)
&  =\left[  i\left(  \eta\right)  \right]  \left(  y_{1}^{d}\wedge\cdots\wedge
y_{p}^{d}\right) \\
&  =\sum_{1\leq\mu\leq p}\left(  -1\right)  ^{\mu-1}\left\langle \eta,y_{\mu
}^{d}\right\rangle y_{1}^{d}\wedge\cdots\wedge\widehat{y_{\mu}^{d}}%
\wedge\cdots\wedge y_{p}^{d},
\end{align}
which can be translated as
\begin{align}
&  i\left(  \left\vert \eta\right\rangle \right)  \left\langle y_{k_{1}%
}\right\vert \wedge\cdots\wedge\left\langle y_{k_{p}}\right\vert \\
&  =\sum_{1\leq\mu\leq p}\left(  -1\right)  ^{\mu-1}\left\langle \left.
y_{k_{\mu}}\right\vert \eta\right\rangle \left\langle y_{k_{1}}\wedge
\cdots\wedge\widehat{y_{k_{\mu}}}\wedge\cdots\wedge y_{k_{p}}\right\vert .
\end{align}

\end{lemma}

\begin{lemma}%
\begin{align}
i\left(  \eta\right)  \left(  z_{1}\cdot\cdots\cdot z_{p}\right)   &  =\left[
i\left(  \eta_{1}\right)  i\left(  \eta_{2}^{d}\right)  \right]  \left(
x_{1}\wedge\cdots\wedge x_{p}\otimes y_{1}^{d}\wedge\cdots\wedge y_{p}%
^{d}\right) \\
&  =\sum_{1\leq\mu,v\leq n}\left(  -1\right)  ^{\mu+v}\left\langle \eta
_{1},y_{\mu}^{d}\right\rangle \left\langle x_{\nu},\eta_{2}^{d}\right\rangle
x_{1}\wedge\cdots\wedge\widehat{x_{\nu}}\wedge\cdots\wedge x_{p}\otimes
y_{1}^{d}\wedge\cdots\wedge\widehat{y_{\mu}^{d}}\wedge\cdots\wedge y_{p}^{d},
\end{align}
where
\begin{equation}
z_{j}=x_{j}\otimes y_{j}^{d};\quad1\leq j\leq n,
\end{equation}
which can be translated as
\begin{align}
&  i\left(  \left\vert x_{1}\right\rangle \left\langle x_{2}\right\vert
\right)  \left(  \left\vert x_{j_{1}}\right\rangle \left\langle x_{k_{1}%
}\right\vert \boxtimes\cdots\boxtimes\left\vert x_{j_{p}}\right\rangle
\left\langle x_{k_{p}}\right\vert \right) \\
&  =\sum_{1\leq v,\mu\leq p}\left(  -1\right)  ^{\mu+v}\left\langle \left.
x_{k_{\mu}}\right\vert x_{1}\right\rangle \left\langle \left.  x_{2}%
\right\vert x_{j\nu}\right\rangle \left\vert x_{_{j_{1}}}\wedge\cdots
\wedge\widehat{x_{j_{\nu}}}\wedge\cdots\wedge x_{j_{p}}\right\rangle
\left\langle x_{k_{1}}\wedge\cdots\wedge\widehat{x_{k_{\mu}}}\wedge
\cdots\wedge x_{k_{p}}\right\vert \\
&  =\sum_{1\leq v,\mu\leq p}\left(  -1\right)  ^{\mu+v}\left\langle \left.
x_{k_{\mu}}\right\vert x_{1}\right\rangle \left\langle \left.  x_{2}%
\right\vert x_{j\nu}\right\rangle \left\vert x_{j_{1}}\right\rangle
\left\langle x_{k_{1}}\right\vert \boxtimes\cdots\boxtimes\left\vert
\widehat{x_{j_{\nu}}}\right\rangle \left\langle x_{k_{\nu}}\right\vert
\boxtimes\cdots\boxtimes\left\vert x_{j_{\mu}}\right\rangle \left\langle
\widehat{x_{k_{\mu}}}\right\vert \boxtimes\cdots\boxtimes\left\vert x_{j_{p}%
}\right\rangle \left\langle x_{k_{p}}\right\vert \\
&  =\sum_{1\leq v,\mu\leq p}\left(  -1\right)  ^{\mu+v}\left\langle \left.
x_{k_{\mu}}\right\vert x_{1}\right\rangle \left\langle \left.  x_{2}%
\right\vert x_{j\nu}\right\rangle \left\vert x_{j_{1}}\right\rangle
\left\langle x_{k_{1}}\right\vert \boxtimes\cdots\boxtimes\left\vert
x_{j_{\mu}}\right\rangle \left\langle x_{k_{\nu}}\right\vert \boxtimes
\cdots\boxtimes\widehat{\left\vert x_{j_{\nu}}\right\rangle \left\langle
x_{k_{\mu}}\right\vert }\boxtimes\cdots\boxtimes\left\vert x_{j_{p}%
}\right\rangle \left\langle x_{k_{p}}\right\vert .
\end{align}

\end{lemma}

The composition product and the geometric product are related by

\begin{proposition}%
\begin{equation}
i\left(  z\right)  \left(  z_{1}\cdot\cdots\cdot z_{p}\right)  =\sum_{1\leq
v\leq p}\left\langle \left.  z\right\vert z_{\nu}\right\rangle z_{1}%
\cdot\cdots\cdot\widehat{z_{\nu}}\cdot\cdots\cdot z_{p}-\sum_{1\leq v<\mu\leq
p}\left(  z_{\mu}\circ z\circ z_{\nu}+z_{\nu}\circ z\circ z_{\mu}\right)
\cdot z_{1}\cdot\cdots\cdot\widehat{z_{\nu}}\cdot\cdots\cdot\widehat{z_{\mu}%
}\cdot\cdots\cdot z_{p}%
\end{equation}
which can be translated as
\begin{align}
&  i\left(  \left\vert \eta_{1}\right\rangle \left\langle \eta_{2}\right\vert
\right)  \left(  \left\vert x_{j_{1}}\right\rangle \left\langle x_{k_{1}%
}\right\vert \boxtimes\cdots\boxtimes\left\vert x_{j_{p}}\right\rangle
\left\langle x_{k_{p}}\right\vert \right)  =\sum_{1\leq v\leq p}\left\langle
\left.  x_{k_{\nu}}\right\vert \eta_{1}\right\rangle \left\langle \left.
\eta_{2}\right\vert x_{j\nu}\right\rangle \left\vert x_{j_{1}}\right\rangle
\left\langle x_{k_{1}}\right\vert \boxtimes\cdots\boxtimes\widehat{\left\vert
x_{j_{\nu}}\right\rangle \left\langle x_{k_{\nu}}\right\vert }\boxtimes
\cdots\boxtimes\left\vert x_{j_{p}}\right\rangle \left\langle x_{k_{p}%
}\right\vert -\\
&  \sum_{1\leq v<\mu\leq p}\left(  \left\vert x_{j\mu}\right\rangle
\left\langle x_{k_{\mu}}\right\vert \circ\left\vert \eta_{1}\right\rangle
\left\langle \eta_{2}\right\vert \circ\left\vert x_{j\nu}\right\rangle
\left\langle x_{k_{\nu}}\right\vert +\left\vert x_{j_{\nu}}\right\rangle
\left\langle x_{k_{\mu}}\right\vert \circ\left\vert \eta_{1}\right\rangle
\left\langle \eta_{2}\right\vert \circ\left\vert x_{j_{\mu}}\right\rangle
\left\langle x_{k_{\mu}}\right\vert \right)  \boxtimes\\
&  \qquad\qquad\qquad\qquad\qquad\qquad\qquad\qquad\qquad\qquad\qquad
\qquad\qquad\boxtimes z_{1}\boxtimes\cdots\boxtimes\widehat{\left\vert
x_{j_{\nu}}\right\rangle \left\langle x_{k_{\nu}}\right\vert }\boxtimes
\cdots\boxtimes\widehat{\left\vert x_{j_{\mu}}\right\rangle \left\langle
x_{k_{\mu}}\right\vert }\boxtimes\cdots\boxtimes z_{p}.
\end{align}

\end{proposition}

\begin{corollary}
Let $z,w\in%
%TCIMACRO{\tbigwedge }%
%BeginExpansion
{\textstyle\bigwedge}
%EndExpansion
\left(  \mathcal{H}^{1},\mathcal{H}^{1d}\right)  $ then
\begin{equation}
i\left(  z\right)  w^{p\cdot}=\left\langle \left.  z\right\vert w\right\rangle
w^{\left(  p-1\right)  \cdot}-\left(  w\circ z\circ w\right)  \cdot w^{\left(
p-2\right)  \cdot}.
\end{equation}

\end{corollary}

\begin{enumerate}
\item
\begin{equation}
i\left(  z\right)  w^{p\cdot}=\left\langle \left.  z\right\vert w\right\rangle
w^{\left(  p-1\right)  \cdot}-\left(  w\circ z\circ w\right)  \cdot w^{\left(
p-2\right)  \cdot}.
\end{equation}


\item
\begin{gather}
\Xi_{\mathcal{H}_{\mathcal{F}}\mathcal{H}_{\mathcal{F}}^{d}}\left(  w\right)
=\left\vert x_{1}\wedge x_{2}\right\rangle \left\langle x_{1}\wedge
x_{2}\right\vert +\left\vert x_{3}\wedge x_{4}\right\rangle \left\langle
x_{3}\wedge x_{4}\right\vert \\
\Xi_{\mathcal{H}_{\mathcal{F}}\mathcal{H}_{\mathcal{F}}^{d}}\left(  z\right)
=\left\vert x_{1}\wedge x_{2}\right\rangle \left\langle x_{1}\wedge
x_{2}\right\vert +\left\vert x_{3}\wedge x_{4}\right\rangle \left\langle
x_{3}\wedge x_{4}\right\vert \\
p=2
\end{gather}


\item
\begin{align}
&  \left(  \left\vert x_{1}\wedge x_{2}\right\rangle \left\langle x_{1}\wedge
x_{2}\right\vert +\left\vert x_{3}\wedge x_{4}\right\rangle \left\langle
x_{3}\wedge x_{4}\right\vert \right)  \left(  \left\vert x_{1}\wedge
x_{2}\right\rangle \left\langle x_{1}\wedge x_{2}\right\vert +\left\vert
x_{3}\wedge x_{4}\right\rangle \left\langle x_{3}\wedge x_{4}\right\vert
\right) \\
&  =\frac{1}{2}\left(  \left\vert x_{1}\wedge x_{2}\right\rangle \left\langle
x_{1}\wedge x_{2}\right\vert +\left\vert x_{3}\wedge x_{4}\right\rangle
\left\langle x_{3}\wedge x_{4}\right\vert \right)
\end{align}


\item
\begin{align}
&  \left\langle \left.  \left(  \left\vert x_{1}\wedge x_{2}\right\rangle
\left\langle x_{1}\wedge x_{2}\right\vert +\left\vert x_{3}\wedge
x_{4}\right\rangle \left\langle x_{3}\wedge x_{4}\right\vert \right)
\right\vert \left(  \left\vert x_{1}\wedge x_{2}\right\rangle \left\langle
x_{1}\wedge x_{2}\right\vert +\left\vert x_{3}\wedge x_{4}\right\rangle
\left\langle x_{3}\wedge x_{4}\right\vert \right)  \right\rangle \\
&  =\left\langle \left.  \overset{}{\left\vert x_{1}\wedge x_{2}\right\rangle
\left\langle x_{1}\wedge x_{2}\right\vert }\right\vert \overset{}{\left\vert
x_{1}\wedge x_{2}\right\rangle \left\langle x_{1}\wedge x_{2}\right\vert
}\right\rangle +\left\langle \left.  \overset{}{\left\vert x_{3}\wedge
x_{4}\right\rangle \left\langle x_{3}\wedge x_{4}\right\vert }\right\vert
\overset{}{\left\vert x_{3}\wedge x_{4}\right\rangle \left\langle x_{3}\wedge
x_{4}\right\vert }\right\rangle \\
&  =\left\langle x_{1}\wedge x_{2}\right\vert \left\vert x_{1}\wedge
x_{2}\right\rangle ^{2}+\left\langle x_{3}\wedge x_{4}\right\vert \left\vert
x_{3}\wedge x_{4}\right\rangle ^{2}=\frac{1}{2}%
\end{align}


\item
\begin{align}
&  i\left(  \left\vert \overset{}{x_{1}\wedge x_{2}}\right\rangle \left\langle
\overset{}{x_{1}\wedge x_{2}}\right\vert +\left\vert \overset{}{x_{3}\wedge
x_{4}}\right\rangle \left\langle \overset{}{x_{3}\wedge x_{4}}\right\vert
\right)  \left(  \left\vert \overset{}{x_{1}\wedge x_{2}}\right\rangle
\left\langle \overset{}{x_{1}\wedge x_{2}}\right\vert +\left\vert
\overset{}{x_{3}\wedge x_{4}}\right\rangle \left\langle \overset{}{x_{3}\wedge
x_{4}}\right\vert \right)  ^{2\cdot}\\
=  &  2i\left(  \left\vert \overset{}{x_{1}\wedge x_{2}}\right\rangle
\left\langle \overset{}{x_{1}\wedge x_{2}}\right\vert +\left\vert
\overset{}{x_{3}\wedge x_{4}}\right\rangle \left\langle \overset{}{x_{3}\wedge
x_{4}}\right\vert \right)  \left(  \left\vert \overset{}{x_{1}\wedge
x_{2}\wedge x_{3}\wedge x_{4}}\right\rangle \left\langle \overset{}{x_{1}%
\wedge x_{2}\wedge x_{3}\wedge x_{4}}\right\vert \right) \\
&  =2\left(
\begin{array}
[c]{c}%
\left\vert i\left(  \left\vert x_{1}\wedge x_{2}\right\rangle \right)
\overset{}{x_{1}\wedge x_{2}\wedge x_{3}\wedge x_{4}}\right\rangle
\left\langle i\left(  \left\vert x_{1}\wedge x_{2}\right\rangle \right)
\overset{}{x_{1}\wedge x_{2}\wedge x_{3}\wedge x_{4}}\right\vert \\
+\left\vert i\left(  \left\vert x_{3}\wedge x_{4}\right\rangle \right)
\overset{}{x_{1}\wedge x_{2}\wedge x_{3}\wedge x_{4}}\right\rangle
\left\langle i\left(  \left\vert x_{3}\wedge x_{4}\right\rangle \right)
\overset{}{x_{1}\wedge x_{2}\wedge x_{3}\wedge x_{4}}\right\vert
\end{array}
\right) \\
&  =2\left(  \left\vert \overset{}{x_{3}\wedge x_{4}}\right\rangle
\left\langle \overset{}{x_{3}\wedge x_{4}}\right\vert +\left\vert
\overset{}{x_{1}\wedge x_{2}}\right\rangle \left\langle \overset{}{x_{1}\wedge
x_{2}}\right\vert \right)  .
\end{align}

\end{enumerate}

\subsection{Poincare Duality}

\subsubsection{Unit Tensors}

The unit tensors in
\begin{equation}
\left\{  \left.
%TCIMACRO{\tbigwedge \nolimits_{p}^{p}}%
%BeginExpansion
{\textstyle\bigwedge\nolimits_{p}^{p}}
%EndExpansion
\left(  \mathcal{H},\mathcal{H}^{d}\right)  \right\vert 0\leq p\leq n\right\}
\cong\left\{  \left.  \Delta_{n}\left(  \mathcal{H}^{1}\right)  \right\vert
0\leq p\leq n\right\}  =\left\{  \left.  \mathcal{B}\left(  \mathcal{H}%
^{p}\right)  \right\vert 0\leq p\leq n\right\}
\end{equation}
are given by
\begin{equation}
\Xi_{\mathcal{H}_{\mathcal{F}}\mathcal{H}_{\mathcal{F}}^{d}}\left(
I_{\wedge\left(  \mathcal{H}^{1}\right)  \otimes\wedge\left(  \mathcal{H}%
^{d}\right)  }\right)  =\left.  P_{\mathcal{H}^{p}}\right\vert _{\mathcal{H}%
^{p}}=I_{\mathcal{H}^{p}},
\end{equation}
where
\begin{gather}
P_{\mathcal{H}^{p}}=\sum_{1\leq j_{1}<\cdots<j_{p}\leq n}\left\vert x_{j_{1}%
}\cdots x_{j_{p}}\right\rangle \left\langle x_{j_{1}}\cdots x_{j_{p}%
}\right\vert =\sum_{1\leq j_{1},\cdots,j_{p}\leq n}\left\vert x_{j_{1}%
}\right\rangle \left\langle x_{j_{1}}\right\vert \cdot\cdots\cdot\left\vert
x_{j_{p}}\right\rangle \left\langle x_{j_{p}}\right\vert =\sum_{1\leq
j_{1},\cdots,j_{p}\leq n}\left\vert x_{j_{1}}\right\rangle \left\langle
x_{j_{1}}\right\vert \boxtimes\cdots\boxtimes\left\vert x_{j_{p}}\right\rangle
\left\langle x_{j_{p}}\right\vert \\
=\left(  P_{\mathcal{H}^{1}}\right)  ^{p\cdot}=\left(  P_{\mathcal{H}^{1}%
}\right)  ^{p\boxtimes}%
\end{gather}
i.e.
\begin{equation}
P_{\mathcal{H}^{p}}=%
%TCIMACRO{\tbigwedge \nolimits^{p}}%
%BeginExpansion
{\textstyle\bigwedge\nolimits^{p}}
%EndExpansion
P_{\mathcal{H}^{1}}=\Xi_{\mathcal{H}_{\mathcal{F}}\mathcal{H}_{\mathcal{F}%
}^{d}}\left(  P_{\mathcal{H}^{1}}\cdot\cdots\cdot P_{\mathcal{H}^{1}}\right)
=P_{\mathcal{H}^{1}}\boxtimes\cdots\boxtimes P_{\mathcal{H}^{1}}.
\end{equation}
Hence
\begin{equation}
I_{\wedge\left(  \mathcal{H}^{1}\right)  \otimes\wedge\left(  \mathcal{H}%
^{d}\right)  }\rightarrow I_{\mathcal{H}_{\mathcal{F}}}=\sum_{0\leq p\leq
n}P_{\mathcal{H}^{p}}.
\end{equation}


\begin{lemma}%
\begin{align}
\left\vert x_{j_{1}}\cdots x_{j_{p}}\right\rangle  &  =\sqrt{p!}\left\vert
x_{j_{1}}\wedge\cdots\wedge x_{j_{p}}\right\rangle \\
&  =\frac{1}{\sqrt{p!}}\left\vert x_{j_{1}}\wedge_{G}\cdots\wedge_{G}x_{j_{p}%
}\right\rangle .
\end{align}

\end{lemma}

\begin{proof}%
\begin{align}
\left\langle \left.  x_{j_{1}}\wedge\cdots\wedge x_{j_{p}}\right\vert
x_{j_{1}}\wedge\cdots\wedge x_{j_{p}}\right\rangle  &  =\left(  \frac{1}%
{p!}\right)  ^{2}\sum_{\lambda,\mu\in S_{p}}\left\langle \left.
x_{j_{\mu\left(  1\right)  }}\otimes\cdots\otimes x_{j_{\mu\left(  p\right)
}}\right\vert x_{j_{\lambda\left(  1\right)  }}\otimes\cdots\otimes
x_{j_{\lambda\left(  p\right)  }}\right\rangle =\left(  \frac{1}{p!}\right)
^{2}p!\,=\frac{1}{p!}\\
&  \Rightarrow\left\langle \left.  x_{j_{1}}\wedge\cdots\wedge x_{j_{p}%
}\right\vert x_{j_{1}}\wedge\cdots\wedge x_{j_{p}}\right\rangle p!\,=1\\
&  \Rightarrow\left\langle \left.  x_{j_{1}}\cdots x_{j_{p}}\right\vert
x_{j_{1}}\cdots x_{j_{p}}\right\rangle \,=1,
\end{align}
thus
\begin{align}
\left\vert x_{j_{1}}\cdots x_{j_{p}}\right\rangle  &  =\sqrt{p!}\left\vert
x_{j_{1}}\wedge\cdots\wedge x_{j_{p}}\right\rangle \\
\frac{1}{\sqrt{p!}}\left\vert x_{j_{1}}\cdots x_{j_{p}}\right\rangle  &
=\left\vert x_{j_{1}}\wedge\cdots\wedge x_{j_{p}}\right\rangle .
\end{align}
The Greub exterior product gives
\begin{align}
&  \left\langle \left.  x_{j_{1}}\wedge_{G}\cdots\wedge_{G}x_{j_{p}%
}\right\vert x_{j_{1}}\wedge_{G}\cdots\wedge_{G}x_{j_{p}}\right\rangle
=\sum_{\lambda,\mu\in S_{p}}\left\langle \left.  x_{j_{\mu\left(  1\right)  }%
}\otimes\cdots\otimes x_{j_{\mu\left(  p\right)  }}\right\vert x_{j_{\lambda
\left(  1\right)  }}\otimes\cdots\otimes x_{j_{\lambda\left(  p\right)  }%
}\right\rangle =p!\,\\
&  \Rightarrow\frac{1}{p!}\left\langle \left.  x_{j_{1}}\wedge_{G}\cdots
\wedge_{G}x_{j_{p}}\right\vert x_{j_{1}}\wedge_{G}\cdots\wedge_{G}x_{j_{p}%
}\right\rangle p!\,=1\Rightarrow\left\langle \left.  x_{j_{1}}\cdots x_{j_{p}%
}\right\vert x_{j_{1}}\cdots x_{j_{p}}\right\rangle \,=1,
\end{align}

\end{proof}

\begin{remark}
The normalized vector is given in terms of the product vector as
\begin{gather}
\left\vert x_{j_{1}}\cdots x_{j_{p}}\right\rangle =\left\Vert \left\vert
x_{j_{1}}\wedge\cdots\wedge x_{j_{p}}\right\rangle \right\Vert ^{-\frac{1}{2}%
}\left\vert x_{j_{1}}\wedge\cdots\wedge x_{j_{p}}\right\rangle =\sqrt
{p!}\left\vert x_{j_{1}}\wedge\cdots\wedge x_{j_{p}}\right\rangle \\
\text{as }\left\Vert \left\vert x_{j_{1}}\wedge\cdots\wedge x_{j_{p}%
}\right\rangle \right\Vert ^{-\frac{1}{2}}=\sqrt{p!}\Leftrightarrow\left\Vert
\left\vert x_{j_{1}}\wedge\cdots\wedge x_{j_{p}}\right\rangle \right\Vert
=\left(  p!\right)  ^{-\frac{1}{2}}.
\end{gather}

\end{remark}

\begin{lemma}%
\begin{equation}
P_{\mathcal{H}^{q}}\cdot P_{\mathcal{H}^{p}}=P_{\mathcal{H}^{p+q}}%
\end{equation}
and%
\begin{equation}
P_{\mathcal{H}^{q}}\cdot_{G}P_{\mathcal{H}^{p}}=\binom{p}{q}P_{\mathcal{H}%
^{p+q}}.
\end{equation}

\end{lemma}

%

\begin{align}
&  \left(  \sum_{1\leq j_{1}<\cdots<j_{p}\leq n}\left\vert x_{j_{1}}\cdots
x_{j_{p}}\right\rangle \left\langle x_{j_{1}}\cdots x_{j_{p}}\right\vert
\right)  \cdot\left(  \sum_{1\leq k_{1}<\cdots<k_{q}\leq n}\left\vert
x_{k_{1}}\cdots x_{k_{q}}\right\rangle \left\langle x_{k_{1}}\cdots x_{k_{q}%
}\right\vert \right) \\
&  =\left(  \frac{1}{p!}\sum_{1\leq j_{1},\cdots,j_{p}\leq n}\left\vert
x_{j_{1}}\cdots x_{j_{p}}\right\rangle \left\langle x_{j_{1}}\cdots x_{j_{p}%
}\right\vert \right)  \cdot\left(  \frac{1}{q!}\sum_{1\leq k_{1},\cdots
,k_{q}\leq n}\left\vert x_{k_{1}}\cdots x_{k_{q}}\right\rangle \left\langle
x_{k_{1}}\cdots x_{k_{q}}\right\vert \right) \\
&  =\left(  \frac{1}{p!}p!\sum_{1\leq j_{1},\cdots,j_{p}\leq n}\left\vert
x_{j_{1}}\wedge\cdots\wedge x_{j_{p}}\right\rangle \left\langle x_{j_{1}%
}\wedge\cdots\wedge x_{j_{p}}\right\vert \right)  \cdot\left(  \frac{1}%
{q!}q!\sum_{1\leq k_{1},\cdots,k_{q}\leq n}\left\vert x_{k_{1}}\wedge
\cdots\wedge x_{k_{q}}\right\rangle \left\langle x_{k_{1}}\wedge\cdots\wedge
x_{k_{q}}\right\vert \right) \\
&  =\sum_{\substack{1\leq j_{1},\cdots,j_{p}\leq n\\1\leq k_{1},\cdots
,k_{q}\leq n}}\left\vert x_{j_{1}}\wedge\cdots\wedge x_{j_{p}}\right\rangle
\left\langle x_{j_{1}}\wedge\cdots\wedge x_{j_{p}}\right\vert \cdot\left\vert
x_{k_{1}}\wedge\cdots\wedge x_{k_{q}}\right\rangle \left\langle x_{k_{1}%
}\wedge\cdots\wedge x_{k_{q}}\right\vert \\
&  =\sum_{\substack{1\leq j_{1},\cdots,j_{p}\leq n\\1\leq k_{1},\cdots
,k_{q}\leq n}}\left\vert x_{j_{1}}\wedge\cdots\wedge x_{j_{p}}\wedge x_{k_{1}%
}\wedge\cdots\wedge x_{k_{q}}\right\rangle \left\langle x_{j_{1}}\wedge
\cdots\wedge x_{j_{p}}\wedge x_{k_{1}}\wedge\cdots\wedge x_{k_{q}}\right\vert
\\
&  =\left(  p+q\right)  !\sum_{1\leq j_{1}<\cdots<j_{p}<k_{1}<\cdots<k_{q}\leq
n}\left\vert x_{j_{1}}\wedge\cdots\wedge x_{j_{p}}\wedge x_{k_{1}}\wedge
\cdots\wedge x_{k_{q}}\right\rangle \left\langle x_{j_{1}}\wedge\cdots\wedge
x_{j_{p}}\wedge x_{k_{1}}\wedge\cdots\wedge x_{k_{q}}\right\vert \\
&  =\sum_{1\leq j_{1}<\cdots<j_{p}<k_{1}<\cdots<k_{q}\leq n}\left\vert
x_{j_{1}}\cdots x_{j_{p}}x_{k_{1}}\cdots x_{k_{q}}\right\rangle \left\langle
x_{j_{1}}\cdots x_{j_{p}}x_{k_{1}}\cdots x_{k_{q}}\right\vert =P_{\mathcal{H}%
^{p+q}}%
\end{align}


\begin{example}%
\begin{align}
\left(  \sum_{1\leq j_{1}\leq n}\left\vert x_{j_{1}}\right\rangle \left\langle
x_{j_{1}}\right\vert \right)  \cdot\left(  \sum_{1\leq k_{1}\leq n}\left\vert
x_{k_{1}}\right\rangle \left\langle x_{k_{1}}\right\vert \right)   &
=\sum_{1\leq j_{1},k_{1}\leq n}\left\vert x_{j_{1}}\wedge x_{k_{1}%
}\right\rangle \left\langle x_{j_{1}}\wedge x_{k_{1}}\right\vert =2\sum_{1\leq
j_{1}<k_{1}\leq n}\left\vert x_{j_{1}}\wedge x_{k_{1}}\right\rangle
\left\langle x_{j_{1}}\wedge x_{k_{1}}\right\vert \\
&  =2\frac{1}{2}\sum_{1\leq j_{1}<k_{1}\leq n}\left\vert x_{j_{1}}x_{k_{1}%
}\right\rangle \left\langle x_{j_{1}}x_{k_{1}}\right\vert =\sum_{1\leq
j_{1}<k_{1}\leq n}\left\vert x_{j_{1}}x_{k_{1}}\right\rangle \left\langle
x_{j_{1}}x_{k_{1}}\right\vert
\end{align}

\end{example}

\begin{example}%
\begin{align}
&  \left(  \sum_{1\leq j_{1}\leq n}\left\vert x_{j_{1}}\right\rangle
\left\langle x_{j_{1}}\right\vert \right)  \cdot\left(  \sum_{1\leq k_{1}\leq
n}\left\vert x_{k_{1}}\right\rangle \left\langle x_{k_{1}}\right\vert \right)
=\left(  \sum_{1\leq j_{1},k_{1}\leq n}\left\vert x_{j_{1}}\wedge x_{k_{1}%
}\right\rangle \left\langle x_{j_{1}}\wedge x_{k_{1}}\right\vert \right)
=2\left(  \sum_{1\leq j_{1}<k_{1}\leq n}\left\vert x_{j_{1}}\wedge x_{k_{1}%
}\right\rangle \left\langle x_{j_{1}}\wedge x_{k_{1}}\right\vert \right) \\
&  =2\frac{1}{2}\left(  \sum_{1\leq j_{1}<k_{1}\leq n}\left\vert x_{j_{1}%
}x_{k_{1}}\right\rangle \left\langle x_{j_{1}}x_{k_{1}}\right\vert \right)
=\sum_{1\leq j_{1}<k_{1}\leq n}\left\vert x_{j_{1}}x_{k_{1}}\right\rangle
\left\langle x_{j_{1}}x_{k_{1}}\right\vert
\end{align}

\end{example}

\begin{example}%
\begin{align}
&  \left(  \sum_{1\leq j_{1}\leq n}\left\vert x_{j_{1}}\right\rangle
\left\langle x_{j_{1}}\right\vert \right)  \cdot\left(  \sum_{1\leq
k_{1},k_{2}\leq n}\left\vert x_{k_{1}}x_{k_{2}}\right\rangle \left\langle
x_{k_{1}}x_{k_{2}}\right\vert \right)  =\left(  \sum_{1\leq j_{1}\leq
n}\left\vert x_{j_{1}}\right\rangle \left\langle x_{j_{1}}\right\vert \right)
\cdot\left(  \sum_{1\leq k_{1},k_{2}\leq n}\left\vert x_{k_{1}}\wedge
x_{k_{2}}\right\rangle \left\langle x_{k_{1}}\wedge x_{k_{2}}\right\vert
\right) \\
&  =\sum_{1\leq j_{1},k_{1},k_{2}\leq n}\left\vert x_{j_{1}}\wedge x_{k_{1}%
}\wedge x_{k_{2}}\right\rangle \left\langle x_{j_{1}}\wedge x_{k_{1}}\wedge
x_{k_{2}}\right\vert =6\sum_{1\leq j_{1}<k_{1}<k_{2}\leq n}\left\vert
x_{j_{1}}\wedge x_{k_{1}}\wedge x_{k_{2}}\right\rangle \left\langle x_{j_{1}%
}\wedge x_{k_{1}}\wedge x_{k_{2}}\right\vert \\
&  =6\frac{1}{6}\sum_{1\leq j_{1}<k_{1}<k_{2}\leq n}\left\vert x_{j_{1}%
}x_{k_{1}}x_{k_{2}}\right\rangle \left\langle x_{j_{1}}x_{k_{1}}x_{k_{2}%
}\right\vert =\sum_{1\leq j_{1}<k_{1}<k_{2}\leq n}\left\vert x_{j_{1}}%
x_{k_{1}}x_{k_{2}}\right\rangle \left\langle x_{j_{1}}x_{k_{1}}x_{k_{2}%
}\right\vert
\end{align}

\end{example}

\begin{lemma}%
\begin{equation}
i\left(  P_{\mathcal{H}^{q}}\right)  P_{\mathcal{H}^{p}}=\binom{n-p+q}%
{q}P_{\mathcal{H}^{p-q}};\quad p\geq q
\end{equation}
\bigskip
\end{lemma}

\begin{proof}%
\begin{equation}
n=4;q=1;p=2
\end{equation}%
\begin{align}
&  i\left(  P_{\mathcal{H}^{1}}\right)  P_{\mathcal{H}^{2}}\\
&  =i\left(  \overset{}{\left\vert x_{1}\right\rangle \left\langle
x_{1}\right\vert +\left\vert x_{2}\right\rangle \left\langle x_{2}\right\vert
+\left\vert x_{3}\right\rangle \left\langle x_{3}\right\vert +\left\vert
x_{4}\right\rangle \left\langle x_{4}\right\vert }\right) \\
&  \qquad\qquad\left(  \overset{}{\left\vert x_{1}x_{2}\right\rangle
\left\langle x_{1}x_{2}\right\vert +\left\vert x_{1}x_{3}\right\rangle
\left\langle x_{1}x_{3}\right\vert +\left\vert x_{1}x_{4}\right\rangle
\left\langle x_{1}x_{4}\right\vert +\left\vert x_{2}x_{3}\right\rangle
\left\langle x_{2}x_{3}\right\vert }+\left\vert x_{2}x_{4}\right\rangle
\left\langle x_{2}x_{4}\right\vert +\left\vert x_{3}x_{4}\right\rangle
\left\langle x_{3}x_{4}\right\vert \right) \\
&  =i\left(  \overset{}{\left\vert x_{1}\right\rangle \left\langle
x_{1}\right\vert }\right)  \left(  \overset{}{\left\vert x_{1}x_{2}%
\right\rangle \left\langle x_{1}x_{2}\right\vert +\left\vert x_{1}%
x_{3}\right\rangle \left\langle x_{1}x_{3}\right\vert +\left\vert x_{1}%
x_{4}\right\rangle \left\langle x_{1}x_{4}\right\vert +\left\vert x_{2}%
x_{3}\right\rangle \left\langle x_{2}x_{3}\right\vert }+\left\vert x_{2}%
x_{4}\right\rangle \left\langle x_{2}x_{4}\right\vert +\left\vert x_{3}%
x_{4}\right\rangle \left\langle x_{3}x_{4}\right\vert \right) \\
&  +i\left(  \overset{}{\left\vert x_{2}\right\rangle \left\langle
x_{2}\right\vert }\right)  \left(  \overset{}{\left\vert x_{1}x_{2}%
\right\rangle \left\langle x_{1}x_{2}\right\vert +\left\vert x_{1}%
x_{3}\right\rangle \left\langle x_{1}x_{3}\right\vert +\left\vert x_{1}%
x_{4}\right\rangle \left\langle x_{1}x_{4}\right\vert +\left\vert x_{2}%
x_{3}\right\rangle \left\langle x_{2}x_{3}\right\vert }+\left\vert x_{2}%
x_{4}\right\rangle \left\langle x_{2}x_{4}\right\vert +\left\vert x_{3}%
x_{4}\right\rangle \left\langle x_{3}x_{4}\right\vert \right) \\
&  +i\left(  \overset{}{\left\vert x_{3}\right\rangle \left\langle
x_{3}\right\vert }\right)  \left(  \overset{}{\left\vert x_{1}x_{2}%
\right\rangle \left\langle x_{1}x_{2}\right\vert +\left\vert x_{1}%
x_{3}\right\rangle \left\langle x_{1}x_{3}\right\vert +\left\vert x_{1}%
x_{4}\right\rangle \left\langle x_{1}x_{4}\right\vert +\left\vert x_{2}%
x_{3}\right\rangle \left\langle x_{2}x_{3}\right\vert }+\left\vert x_{2}%
x_{4}\right\rangle \left\langle x_{2}x_{4}\right\vert +\left\vert x_{3}%
x_{4}\right\rangle \left\langle x_{3}x_{4}\right\vert \right) \\
&  +i\left(  \overset{}{\left\vert x_{4}\right\rangle \left\langle
x_{4}\right\vert }\right)  \left(  \overset{}{\left\vert x_{1}x_{2}%
\right\rangle \left\langle x_{1}x_{2}\right\vert +\left\vert x_{1}%
x_{3}\right\rangle \left\langle x_{1}x_{3}\right\vert +\left\vert x_{1}%
x_{4}\right\rangle \left\langle x_{1}x_{4}\right\vert +\left\vert x_{2}%
x_{3}\right\rangle \left\langle x_{2}x_{3}\right\vert }+\left\vert x_{2}%
x_{4}\right\rangle \left\langle x_{2}x_{4}\right\vert +\left\vert x_{3}%
x_{4}\right\rangle \left\langle x_{3}x_{4}\right\vert \right)
\end{align}%
\begin{align}
&  =\left(  \left\vert x_{2}\right\rangle \left\langle x_{2}\right\vert
+\left\vert x_{3}\right\rangle \left\langle x_{3}\right\vert +\left\vert
x_{4}\right\rangle \left\langle x_{4}\right\vert \right)  +\left(  \left\vert
x_{1}\right\rangle \left\langle x_{1}\right\vert +\left\vert x_{3}%
\right\rangle \left\langle x_{3}\right\vert +\left\vert x_{4}\right\rangle
\left\langle x_{4}\right\vert \right) \\
&  +\left(  \left\vert x_{1}\right\rangle \left\langle x_{1}\right\vert
+\left\vert x_{2}\right\rangle \left\langle x_{2}\right\vert +\left\vert
x_{4}\right\rangle \left\langle x_{4}\right\vert \right)  +\left(  \left\vert
x_{1}\right\rangle \left\langle x_{1}\right\vert +\left\vert x_{2}%
\right\rangle \left\langle x_{2}\right\vert +\left\vert x_{3}\right\rangle
\left\langle x_{3}\right\vert \right) \\
&  =3\left(  \left\vert x_{1}\right\rangle \left\langle x_{1}\right\vert
+\left\vert x_{2}\right\rangle \left\langle x_{2}\right\vert +\left\vert
x_{3}\right\rangle \left\langle x_{3}\right\vert +\left\vert x_{4}%
\right\rangle \left\langle x_{4}\right\vert \right)  =3P_{\mathcal{H}^{1}}\\
&  =\binom{4-2+1}{1}P_{\mathcal{H}^{2}}=\binom{3}{1}P_{\mathcal{H}^{2}%
}=3P_{\mathcal{H}^{1}}.
\end{align}%
\begin{equation}
n=4;q=1;p=3
\end{equation}%
\begin{align}
&  i\left(  P_{\mathcal{H}^{1}}\right)  P_{\mathcal{H}^{3}}\\
&  =i\left(  \overset{}{\left\vert x_{1}\right\rangle \left\langle
x_{1}\right\vert +\left\vert x_{2}\right\rangle \left\langle x_{2}\right\vert
+\left\vert x_{3}\right\rangle \left\langle x_{3}\right\vert +\left\vert
x_{4}\right\rangle \left\langle x_{4}\right\vert }\right)  \left(
\overset{}{\left\vert x_{1}x_{2}x_{3}\right\rangle \left\langle x_{1}%
x_{2}x_{3}\right\vert +\left\vert x_{1}x_{2}x_{4}\right\rangle \left\langle
x_{1}x_{2}x_{4}\right\vert +\left\vert x_{1}x_{3}x_{4}\right\rangle
\left\langle x_{1}x_{3}x_{4}\right\vert +\left\vert x_{2}x_{3}x_{4}%
\right\rangle \left\langle x_{2}x_{3}x_{4}\right\vert }\right) \\
&  =i\left(  \overset{}{\left\vert x_{1}\right\rangle \left\langle
x_{1}\right\vert }\right)  \left(  \overset{}{\left\vert x_{1}x_{2}%
x_{3}\right\rangle \left\langle x_{1}x_{2}x_{3}\right\vert +\left\vert
x_{1}x_{2}x_{4}\right\rangle \left\langle x_{1}x_{2}x_{4}\right\vert
+\left\vert x_{1}x_{3}x_{4}\right\rangle \left\langle x_{1}x_{3}%
x_{4}\right\vert +\left\vert x_{2}x_{3}x_{4}\right\rangle \left\langle
x_{2}x_{3}x_{4}\right\vert }\right) \\
&  +i\left(  \overset{}{\left\vert x_{2}\right\rangle \left\langle
x_{2}\right\vert }\right)  \left(  \overset{}{\left\vert x_{1}x_{2}%
x_{3}\right\rangle \left\langle x_{1}x_{2}x_{3}\right\vert +\left\vert
x_{1}x_{2}x_{4}\right\rangle \left\langle x_{1}x_{2}x_{4}\right\vert
+\left\vert x_{1}x_{3}x_{4}\right\rangle \left\langle x_{1}x_{3}%
x_{4}\right\vert +\left\vert x_{2}x_{3}x_{4}\right\rangle \left\langle
x_{2}x_{3}x_{4}\right\vert }\right) \\
&  +i\left(  \overset{}{\left\vert x_{3}\right\rangle \left\langle
x_{3}\right\vert }\right)  \left(  \overset{}{\left\vert x_{1}x_{2}%
x_{3}\right\rangle \left\langle x_{1}x_{2}x_{3}\right\vert +\left\vert
x_{1}x_{2}x_{4}\right\rangle \left\langle x_{1}x_{2}x_{4}\right\vert
+\left\vert x_{1}x_{3}x_{4}\right\rangle \left\langle x_{1}x_{3}%
x_{4}\right\vert +\left\vert x_{2}x_{3}x_{4}\right\rangle \left\langle
x_{2}x_{3}x_{4}\right\vert }\right) \\
&  +i\left(  \overset{}{\left\vert x_{4}\right\rangle \left\langle
x_{4}\right\vert }\right)  \left(  \overset{}{\left\vert x_{1}x_{2}%
x_{3}\right\rangle \left\langle x_{1}x_{2}x_{3}\right\vert +\left\vert
x_{1}x_{2}x_{4}\right\rangle \left\langle x_{1}x_{2}x_{4}\right\vert
+\left\vert x_{1}x_{3}x_{4}\right\rangle \left\langle x_{1}x_{3}%
x_{4}\right\vert +\left\vert x_{2}x_{3}x_{4}\right\rangle \left\langle
x_{2}x_{3}x_{4}\right\vert }\right) \\
&  =\left\vert x_{2}x_{3}\right\rangle \left\langle x_{2}x_{3}\right\vert
+\left\vert x_{2}x_{4}\right\rangle \left\langle x_{2}x_{4}\right\vert
+\left\vert x_{3}x_{4}\right\rangle \left\langle x_{3}x_{4}\right\vert
+\left\vert x_{1}x_{3}\right\rangle \left\langle x_{1}x_{3}\right\vert
+\left\vert x_{1}x_{4}\right\rangle \left\langle x_{1}x_{4}\right\vert
+\left\vert x_{3}x_{4}\right\rangle \left\langle x_{3}x_{4}\right\vert \\
&  +\left\vert x_{1}x_{2}\right\rangle \left\langle x_{1}x_{2}\right\vert
+\left\vert x_{1}x_{4}\right\rangle \left\langle x_{1}x_{4}\right\vert
+\left\vert x_{2}x_{4}\right\rangle \left\langle x_{2}x_{4}\right\vert
+\left\vert x_{1}x_{2}\right\rangle \left\langle x_{1}x_{2}\right\vert
+\left\vert x_{1}x_{3}\right\rangle \left\langle x_{1}x_{3}\right\vert
+\left\vert x_{2}x_{3}\right\rangle \left\langle x_{2}x_{3}\right\vert \\
&  =2\left(  \left\vert x_{2}x_{3}\right\rangle \left\langle x_{2}%
x_{3}\right\vert +\left\vert x_{2}x_{4}\right\rangle \left\langle x_{2}%
x_{4}\right\vert +\left\vert x_{3}x_{4}\right\rangle \left\langle x_{3}%
x_{4}\right\vert +\left\vert x_{1}x_{3}\right\rangle \left\langle x_{1}%
x_{3}\right\vert +\left\vert x_{1}x_{4}\right\rangle \left\langle x_{1}%
x_{4}\right\vert +\left\vert x_{3}x_{4}\right\rangle \left\langle x_{3}%
x_{4}\right\vert +\left\vert x_{1}x_{2}\right\rangle \left\langle x_{1}%
x_{2}\right\vert \right) \\
&  =\binom{4-3+1}{1}P_{\mathcal{H}^{3}}=\binom{2}{1}P_{\mathcal{H}^{3}%
}=2P_{\mathcal{H}^{3}}.
\end{align}%
\begin{equation}
q=1
\end{equation}%
\begin{align}
&  i\left(  P_{\mathcal{H}^{1}}\right)  P_{\mathcal{H}^{p}}=i\left(
\sum_{1\leq j\leq n}\left\vert x_{j}\right\rangle \left\langle x_{j}%
\right\vert \right)  \left(  \sum_{1\leq j_{1},\cdots,j_{p}\leq n}\left\vert
x_{j_{1}}\right\rangle \left\langle x_{j_{1}}\right\vert \wedge\cdots
\wedge\left\vert x_{j_{p}}\right\rangle \left\langle x_{j_{p}}\right\vert
\right) \\
&  =i\left(  \sum_{1\leq j\leq n}\left\vert x_{j}\right\rangle \left\langle
x_{j}\right\vert \right)  \left(  \sum_{1\leq j_{1}<\cdots<j_{p}\leq
n}\left\vert x_{j_{1}}\cdots x_{j_{p}}\right\rangle \left\langle x_{j_{1}%
}\cdots x_{j_{p}}\right\vert \right) \\
&  =\sum_{1\leq j\leq n}\sum_{1\leq j_{1}<\cdots<j_{p}\leq n}\left\vert
i\left(  \left\langle x_{j}\right\vert \right)  \left(  x_{j_{1}}\cdots
x_{j_{p}}\right)  \right\rangle \left\langle i\left(  \left\langle
x_{j}\right\vert \right)  \left(  x_{j_{1}}\cdots x_{j_{p}}\right)
\right\vert \\
&  =\sum_{1\leq j\leq n}\sum_{1\leq\nu\leq p}\sum_{1\leq j_{1}<\cdots
<j_{p}\leq n}\left\langle \left.  x_{j}\right\vert x_{j_{\nu\left(  1\right)
}}\right\rangle \left\langle \left.  x_{j_{\nu\left(  1\right)  }}\right\vert
x_{j}\right\rangle p!\left\vert x_{j_{1}}\wedge\cdots\wedge\widehat{x_{j_{\nu
\left(  1\right)  }}}\wedge\cdots\wedge x_{j_{p}}\right\rangle \left\langle
x_{j_{1}}\wedge\cdots\wedge\widehat{x_{j_{\nu\left(  1\right)  }}}\wedge
\cdots\wedge x_{j_{p}}\right\vert \\
&  =\sum_{1\leq j\leq n}\sum_{1\leq\nu\leq p}\sum_{1\leq j_{1}<\cdots
<j_{p}\leq n}\left\langle \left.  x_{j}\right\vert x_{j_{\nu\left(  1\right)
}}\right\rangle \left\langle \left.  x_{j_{\nu\left(  1\right)  }}\right\vert
x_{j}\right\rangle \frac{p!}{\left(  p-1\right)  !}\left\vert x_{j_{1}}%
\cdots\widehat{x_{j_{\nu\left(  1\right)  }}}\cdots x_{j_{p}}\right\rangle
\left\langle x_{j_{1}}\cdots\widehat{x_{j_{\nu\left(  1\right)  }}}\cdots
x_{j_{p}}\right\vert \\
&  =\sum_{1\leq j\leq n}\sum_{1\leq\nu\leq p}\sum_{1\leq j_{1}<\cdots
<j_{p}\leq n}\left\langle \left.  x_{j}\right\vert x_{j_{\nu\left(  1\right)
}}\right\rangle \left\langle \left.  x_{j_{\nu\left(  1\right)  }}\right\vert
x_{j}\right\rangle p\left\vert x_{j_{1}}\cdots\widehat{x_{j_{\nu\left(
1\right)  }}}\cdots x_{j_{p}}\right\rangle \left\langle x_{j_{1}}%
\cdots\widehat{x_{j_{\nu\left(  1\right)  }}}\cdots x_{j_{p}}\right\vert \\
&  =\sum_{1\leq j_{1}<\cdots<j_{p}\leq n}\binom{n}{p}p\left\vert x_{j_{1}%
}\cdots\widehat{x_{j_{\nu\left(  1\right)  }}}\cdots x_{j_{p}}\right\rangle
\left\langle x_{j_{1}}\cdots\widehat{x_{j_{\nu\left(  1\right)  }}}\cdots
x_{j_{p}}\right\vert =\sum_{1\leq j_{1}<\cdots<j_{p}\leq n}\frac{n\left(
n-1\right)  \cdots\left(  n-p\right)  }{\left(  p-1\right)  !}\left\vert
x_{j_{1}}\cdots\widehat{x_{j_{\nu\left(  1\right)  }}}\cdots x_{j_{p}%
}\right\rangle \left\langle x_{j_{1}}\cdots\widehat{x_{j_{\nu\left(  1\right)
}}}\cdots x_{j_{p}}\right\vert \\
&  =\sum_{1\leq j_{1}<\cdots<j_{p}\leq n}\frac{n\left(  n-1\right)
\cdots\left(  n-p\right)  }{\left(  p-1\right)  !}\left\vert x_{j_{1}}%
\cdots\widehat{x_{j_{\nu\left(  1\right)  }}}\cdots x_{j_{p}}\right\rangle
\left\langle x_{j_{1}}\cdots\widehat{x_{j_{\nu\left(  1\right)  }}}\cdots
x_{j_{p}}\right\vert
\end{align}

\end{proof}

%

\begin{align}
\left\langle \left.  i\left(  P_{\mathcal{H}^{q}}\right)  P_{\mathcal{H}^{p}%
}\right\vert P_{\mathcal{H}^{p-q}}\right\rangle  &  =\operatorname{Tr}\left\{
\left(  i\left(  P_{\mathcal{H}^{q}}\right)  P_{\mathcal{H}^{p}}\right)
^{\dagger}P_{\mathcal{H}^{p-q}}\right\}  =\left\langle \left.  \binom
{n-p+q}{q}P_{\mathcal{H}^{p-q}}\right\vert P_{\mathcal{H}^{p-q}}\right\rangle
\\
&  =\binom{n-p+q}{q}\operatorname{Tr}\left\{  P_{\mathcal{H}^{p-q}}\right\}
=\binom{n-p+q}{q}\binom{n}{p-q}%
\end{align}%
\begin{equation}
\binom{n-p+q}{q}\binom{n}{p-q}=\frac{\left(  n-p+q\right)  !}{\left(
n-p\right)  !\left(  q\right)  !}\frac{\left(  n\right)  !}{\left(
p-q\right)  !\left(  n-p+q\right)  !}=\frac{\left(  n\right)  !}{\left(
n-p\right)  !\left(  p-q\right)  !\left(  q\right)  !}%
\end{equation}%
\begin{align}
n  &  =6;p=3;q=1\\
&  \left(  6-3\right)  !\left(  2\right)  !\left(  1\right)  !
\end{align}%
\begin{align}
\left\langle \left.  i\left(  P_{\mathcal{H}^{q}}\right)  P_{\mathcal{H}^{p}%
}\right\vert P_{\mathcal{H}^{p-q}}\right\rangle  &  =\left\langle \left.
P_{\mathcal{H}^{p}}\right\vert \mu\left(  P_{\mathcal{H}^{q}}\right)
P_{\mathcal{H}^{p-q}}\right\rangle =\left\langle \left.  P_{\mathcal{H}^{p}%
}\right\vert P_{\mathcal{H}^{q}}\boxtimes P_{\mathcal{H}^{p-q}}\right\rangle
=\alpha\binom{n}{p}=\alpha\frac{\left(  n\right)  !}{\left(  p\right)
!\left(  n-p\right)  !}=\frac{\left(  n\right)  !}{\left(  n-p\right)
!\left(  p-q\right)  !\left(  q\right)  !}\\
\alpha &  =\frac{\left(  p\right)  !}{\left(  p-q\right)  !\left(  q\right)
!}=\binom{p}{q}%
\end{align}%
\begin{equation}
P_{\mathcal{H}^{q}}\boxtimes P_{\mathcal{H}^{p-q}}=\binom{p}{q}P_{\mathcal{H}%
^{p}}%
\end{equation}


\begin{definition}
In finite dimensions define the linear maps%
\begin{align}
D_{e} &  :%
%TCIMACRO{\tbigwedge }%
%BeginExpansion
{\textstyle\bigwedge}
%EndExpansion
\left(  \mathcal{H}\right)  \rightarrow%
%TCIMACRO{\tbigwedge }%
%BeginExpansion
{\textstyle\bigwedge}
%EndExpansion
\left(  \mathcal{H}^{d}\right)  \\
D^{e} &  :%
%TCIMACRO{\tbigwedge }%
%BeginExpansion
{\textstyle\bigwedge}
%EndExpansion
\left(  \mathcal{H}^{d}\right)  \rightarrow%
%TCIMACRO{\tbigwedge }%
%BeginExpansion
{\textstyle\bigwedge}
%EndExpansion
\left(  \mathcal{H}\right)
\end{align}
by%
\begin{align}
D_{e}u &  =i\left(  u\right)  x^{sd};\quad u\in%
%TCIMACRO{\tbigwedge }%
%BeginExpansion
{\textstyle\bigwedge}
%EndExpansion
\left(  \mathcal{H}\right)  \\
D^{e}u &  =i\left(  u^{d}\right)  x^{s};\quad u^{d}\in%
%TCIMACRO{\tbigwedge }%
%BeginExpansion
{\textstyle\bigwedge}
%EndExpansion
\left(  \mathcal{H}^{d}\right)  ,
\end{align}
where the psuedo scalar is given by
\begin{align}
x^{s} &  =x_{1}\wedge\cdots\wedge x_{n}\\
x^{sd} &  =x_{1}^{d}\wedge\cdots\wedge x_{n}^{d}.
\end{align}
$\lceil$%
\begin{align}
D_{e}\left(  x_{1}\right)   &  =i\left(  \left\vert x_{1}\right\rangle
\right)  \left\langle x_{1}\wedge\cdots\wedge x_{n}\right\vert \\
D^{e}\left(  x_{1}\right)   &  =i\left(  \left\langle x_{1}\right\vert
\right)  \left\vert x_{1}\wedge\cdots\wedge x_{n}\right\rangle =\left\vert
x_{2}\wedge\cdots\wedge x_{n}\right\rangle
\end{align}
$\rfloor$
\end{definition}

The psuedo scalar is normalized wrt the scalar and inner product
\begin{equation}
\left\langle x^{d},x^{sd}\right\rangle =\left\langle \left.  x^{s}\right\vert
x^{s}\right\rangle =1.
\end{equation}
The projector onto the psuedo scalar is given by
\begin{equation}
x^{s}\otimes x^{sd}=\left\vert x^{s}\right\rangle \left\langle x^{s}%
\right\vert =\mathcal{I}_{n}=P_{\mathcal{H}^{n}}%
\end{equation}


The $D_{e}$ has the properties

\begin{enumerate}
\item
\begin{align}
D_{e}\left(  1\right)   &  =x^{sd}=\left\langle x^{s}\right\vert =D_{e}\left(
\left\vert \phi^{s}\right\rangle \right) \\
D^{e}\left(  1^{d}\right)   &  =x^{s}=\left\vert x^{s}\right\rangle
=D^{e}\left(  \left\langle \phi^{sd}\right\vert \right)
\end{align}


\item
\begin{align}
D_{\lambda e}  &  =\lambda^{-1}D_{e}\\
D^{\lambda e}  &  =\lambda D^{e}%
\end{align}


\item
\begin{align}
D_{e}\left(  u\wedge v\right)   &  =i\left(  v\right)  D_{e}\left(  u\right)
;\quad u,v\in%
%TCIMACRO{\tbigwedge }%
%BeginExpansion
{\textstyle\bigwedge}
%EndExpansion
\left(  \mathcal{H}\right) \\
D_{e}\left(  u^{d}\wedge v^{d}\right)   &  =i\left(  v^{d}\right)
D_{e}\left(  u^{d}\right)  ;\quad u^{d},v^{d}\in%
%TCIMACRO{\tbigwedge }%
%BeginExpansion
{\textstyle\bigwedge}
%EndExpansion
\left(  \mathcal{H}^{d}\right)
\end{align}


\item
\begin{align}
D_{p}  &  :%
%TCIMACRO{\tbigwedge \nolimits^{p}}%
%BeginExpansion
{\textstyle\bigwedge\nolimits^{p}}
%EndExpansion
\left(  \mathcal{H}\right)  \rightarrow%
%TCIMACRO{\tbigwedge \nolimits^{n-p}}%
%BeginExpansion
{\textstyle\bigwedge\nolimits^{n-p}}
%EndExpansion
\left(  \mathcal{H}^{d}\right)  ;\quad0\leq p\leq n\\
D^{p}  &  :%
%TCIMACRO{\tbigwedge \nolimits^{p}}%
%BeginExpansion
{\textstyle\bigwedge\nolimits^{p}}
%EndExpansion
\left(  \mathcal{H}^{d}\right)  \rightarrow%
%TCIMACRO{\tbigwedge \nolimits^{n-p}}%
%BeginExpansion
{\textstyle\bigwedge\nolimits^{n-p}}
%EndExpansion
\left(  \mathcal{H}\right)  ;\quad0\leq p\leq n,
\end{align}
where
\begin{equation}
D_{p}=\left.  D_{e}\right\vert _{\mathcal{H}^{P}}%
\end{equation}
and
\begin{equation}
D^{p}=\left.  D^{e}\right\vert _{\mathcal{H}^{P}}.
\end{equation}
$\lceil$%
\begin{equation}
D_{p}\left(  \left\vert x_{1}\wedge\cdots\wedge x_{p}\right\rangle \right)
=D_{e}\left(  \left\vert x_{1}\wedge\cdots\wedge x_{p}\right\rangle \right)
=i\left(  \left\langle x_{1}\wedge\cdots\wedge x_{p}\right\vert \right)
\left\vert x_{1}\wedge\cdots\wedge x_{n}\right\rangle =\left\vert
x_{p+1}\wedge\cdots\wedge x_{n}\right\rangle
\end{equation}
$\rfloor$
\end{enumerate}

\begin{theorem}
$D_{e}$ and $D^{e}$ preserve the scalar products

\begin{enumerate}
\item
\begin{equation}
\left\langle D_{e}u,D^{e}v^{d}\right\rangle =\left\langle u,v^{d}\right\rangle
=\left\langle \left.  D_{e}v\right\vert D_{e}u\right\rangle =\left\langle
\left.  v\right\vert u\right\rangle
\end{equation}


\item The duals $D_{p}^{d}$ and $\left(  D^{p}\right)  ^{d}$ are given by
\begin{equation}
D_{p}^{d}=\left(  -1\right)  ^{p\left(  n-p\right)  }D_{\left(  n-p\right)
};\quad0\leq p\leq n
\end{equation}
and
\begin{equation}
D^{pd}=\left(  -1\right)  ^{p\left(  n-p\right)  }D^{\left(  n-p\right)
};\quad0\leq p\leq n
\end{equation}


\item
\begin{equation}
D^{p}\circ D_{\left(  n-p\right)  }=\left(  -1\right)  ^{p\left(  n-p\right)
}I_{n};\quad0\leq p\leq n
\end{equation}
and
\begin{equation}
D^{p}\circ D_{\left(  n-p\right)  }=\left(  -1\right)  ^{p\left(  n-p\right)
}I_{n};\quad0\leq p\leq n
\end{equation}
In particular $D_{e}$,$D^{e},D_{e}$ and $D^{e}$ are linear isomorphisms.
\end{enumerate}
\end{theorem}

In the case of Hilbert spaces

\begin{enumerate}
\item
\begin{equation}
D_{e}u=i\left(  u\right)  \left\vert x_{s}\right\rangle ;\left\vert
x_{s}\right\rangle \in\mathcal{H}^{p}%
\end{equation}


\item
\begin{equation}
D_{e}:\mathcal{H}^{p}\rightarrow\mathcal{H}^{n-p}%
\end{equation}


\item
\begin{equation}
\left\langle \left.  D_{e}u\right\vert D_{e}v\right\rangle _{\mathcal{H}%
_{\mathcal{F}}}=\left\langle \left.  u\right\vert v\right\rangle
_{\mathcal{H}_{\mathcal{F}}};u,v\in\mathcal{H}_{\mathcal{F}}%
\end{equation}


\item
\begin{equation}
D_{p}^{\dagger}=\left(  -1\right)  ^{p\left(  n-p\right)  }D_{n-p}%
\end{equation}


\item If $n=2m$ then
\begin{equation}
D_{m}:\mathcal{H}^{p}\rightarrow\mathcal{H}^{p}%
\end{equation}
and%
\begin{equation}
D_{m}^{\dagger}=\left(  -1\right)  ^{m}D_{m}%
\end{equation}

\end{enumerate}

\subsubsection{Intersection Product "$\cap"$}%

\begin{equation}
u\cap v=D^{e}\left[  \left(  D^{e}\right)  ^{-1}u\wedge\left(  D^{e}\right)
^{-1}v\right]  ;u,v\in%
%TCIMACRO{\tbigwedge }%
%BeginExpansion
{\textstyle\bigwedge}
%EndExpansion
\left(  \mathcal{H}^{1}\right)  .
\end{equation}
If $u\in\mathcal{H}^{p}$ and $v\in\mathcal{H}^{q}$ then $u\cap v\in
\mathcal{H}^{p+q-n}.$ In particular $u\cap v=0$ if $p+q>2n$ or $p+q<n.$ It is
immediate from the definition that
\begin{equation}
u\cap v=\left(  -1\right)  ^{\left(  n-p\right)  \left(  n-q\right)  }v\cap
u;u\in\mathcal{H}^{p},v\in\mathcal{H}^{q}.
\end{equation}
The unit element in the \emph{Intersection Algebra} is $x_{s}$ i.e.
\begin{equation}
u\cap x_{s}=x_{s}\cap u.
\end{equation}
The map $D_{e}$ is an isomorphism from the intersection algebra to the
exterior algebra i.e.
\begin{equation}
D_{e}\left(  u\cap v\right)  =D_{e}\left(  u\right)  \wedge D_{e}\left(
v\right)  ;u,v\in%
%TCIMACRO{\tbigwedge }%
%BeginExpansion
{\textstyle\bigwedge}
%EndExpansion
\left(  \mathcal{H}^{1}\right)  .
\end{equation}

\end{document}